\section{Proofs for Section~\ref{sec:rates}}
\label{app:rates}

This section proves the results of Section~\ref{sec:rates}: the margin condition and its examples, the rate of the
split-sample bound under a margin, the margins of touching contaminations, the excess of the population ramp, the
parametric rate, up to logarithmic factors, with the Laplace kernel, and the two-point lower bound and the minimax
corollary. It also states and proves the converse for the analytic baseline.

\subsection{Proof of Example~\ref{ex:separated}}
\label{sec:exproof}

\begin{proof}
Every $\alpha$-trimming keeps enough mass where $f=\pm1$ for $f$ to separate the trimmed means. Since
$P-Q=\alpha^\star(P_0-Q_0)$ and $P_0$, $Q_0$ are concentrated on the disjoint sets $S_P$, $S_Q$,
$\TV(P,Q)=\alpha^\star\TV(P_0,Q_0)=\alpha^\star$. Fix $\alpha\in[0,\alpha^\star]$ and an $\alpha$-trimming $P'=gP$, with
$0\le g\le\beta$. The law $P$ is concentrated on $S_R\cup S_P$ and gives $S_R$ the mass $1-\alpha^\star$, and $f$ is $0$
on $S_R$ and $1$ on $S_P$, so
\[
\E_{P'}f=\int_{S_P}g\,dP=1-\int_{S_R}g\,dP\ge1-\beta(1-\alpha^\star)=\frac{\alpha^\star-\alpha}{1-\alpha}.
\]
In the same way $\E_{Q'}f\le-(\alpha^\star-\alpha)/(1-\alpha)$ for every $Q'\in R_\alpha(Q)$. The function $f$ is not $0$,
since it equals $1$ on $S_P$, and by the reproducing property and the Cauchy--Schwarz inequality
\[
\MMD(P',Q')=\|\mu_{P'}-\mu_{Q'}\|_\Hh\ge\frac{\langle\mu_{P'}-\mu_{Q'},f\rangle_\Hh}{\|f\|_\Hh}
=\frac{\E_{P'}f-\E_{Q'}f}{\|f\|_\Hh}\ge\frac{2(\alpha^\star-\alpha)}{(1-\alpha)\|f\|_\Hh}.
\]
The infimum over the two trimmings gives the first inequality of the example, and $1-\alpha\le1$ the second. At
$\alpha=\alpha^\star-\epsilon$ it reads $D_{\alpha^\star-\epsilon}(P,Q)\ge2\epsilon/\|f\|_\Hh$, which is
Assumption~\ref{ass:margin} with $c_0=2/\|f\|_\Hh$ and $\gamma=1$.
\end{proof}

For finite sets and a strictly positive definite kernel the Gram matrix of their points is invertible, so a combination
$\sum_ic_ik(\cdot,z_i)$ takes any prescribed values on them. For a Mat\'ern kernel on $\R^d$ the RKHS is a Sobolev space
with an equivalent norm \citep[Chapter~10]{wendland2005scattered}, which contains a smooth compactly supported function equal to
$1$ on the $\varepsilon/3$-neighborhood of $S_P$, to $-1$ on that of $S_Q$ and to $0$ outside their
$2\varepsilon/3$-neighborhoods, when $S_P$ and $S_Q$ are bounded and $\varepsilon>0$ is the smallest distance between
two of the three sets.

\subsection{A general margin and the convexity of the path}
\label{sec:generalmargin}

Assumption~\ref{ass:margin} asks for a polynomial lower bound on the population path below $\alpha^\star$. A lower
bound with no speed always exists, and Proposition~\ref{prop:newrates}(i), applied to $(1-\alpha^\star)\omega$, gives
the split-sample bound a rate for it.

\begin{lemma}[A margin always exists]
\label{lem:marginexists}
Let Assumption~\ref{ass:regularity} hold, let $0<\alpha^\star<1$, and for $\epsilon\in[0,\alpha^\star]$ put
\[
m(\epsilon)\coloneqq\min_{\epsilon'\in[\epsilon,\alpha^\star]}D_{\alpha^\star-\epsilon'}(P,Q),
\qquad
\omega(\epsilon)\coloneqq\frac{\epsilon}{\alpha^\star}\,m(\epsilon).
\]
Then $\omega$ is continuous and strictly increasing on $[0,\alpha^\star]$, $\omega(0)=0$, and
$D_{\alpha^\star-\epsilon}(P,Q)\ge\omega(\epsilon)$ for every $\epsilon\in[0,\alpha^\star]$.
\end{lemma}

\begin{proof}
The map $\epsilon'\mapsto D_{\alpha^\star-\epsilon'}$ is continuous on $[0,\alpha^\star]$ by Lemma~\ref{lem:continuity},
zero at $\epsilon'=0$ and positive elsewhere by \eqref{eq:zeroeq}. Its running minimum $m$ is therefore attained,
continuous and nondecreasing, with $m(0)=0$ and $m(\epsilon)>0$ for $\epsilon>0$. Hence $\omega$ is continuous with
$\omega(0)=0$, and for $0<\epsilon_1<\epsilon_2$, $\omega(\epsilon_2)\ge(\epsilon_2/\alpha^\star)m(\epsilon_1)>\omega(\epsilon_1)>0$.
Finally $\omega(\epsilon)\le m(\epsilon)\le D_{\alpha^\star-\epsilon}$, as $\epsilon\le\alpha^\star$.
\end{proof}

The next lemma, stated in words in Section~\ref{sec:rate}, shows which margins are possible at all, through the
convexity of the rescaled path. For $\alpha\in[0,1)$ let
\[
\bar{\mathcal C}_\alpha(P)\coloneqq(1-\alpha)\,\mathcal C_\alpha(P)=\big\{\E_P[h(X)\phi(X)]:\ 0\le h\le1,\ \E_P[h]=1-\alpha\big\},
\]
the embeddings of the parts $P'\le P$ of mass $1-\alpha$, that is, of the $\alpha$-trimmings before they are
renormalized, with $h=(1-\alpha)g$ for $g$ as in \eqref{eq:region}, and likewise $\bar{\mathcal C}_\alpha(Q)$. By
Proposition~\ref{prop:georep}, which holds for any probability measures and so for $P_n$ and $Q_m$, the rescaled path
is $\bar D_\alpha=(1-\alpha)D_\alpha(P,Q)=\dist(\bar{\mathcal C}_\alpha(P),\bar{\mathcal C}_\alpha(Q))$.

\begin{lemma}[Convexity of the path]
\label{lem:pathconvex}
Let Assumption~\ref{ass:regularity} hold, let $P,Q$ be any probability measures on $\mathcal X$, empirical measures
included, and put $\bar D_{\alpha}\coloneqq(1-\alpha)D_\alpha(P,Q)$ for $\alpha\in[0,1)$. Then $\bar D$ is convex and
nonincreasing. If moreover $0<\alpha^\star<1$, then:
\begin{enumerate}[label=(\roman*),leftmargin=*]
\item $\bar D$ is strictly decreasing on $[0,\alpha^\star]$, and for all $0\le\alpha_1<\alpha_2<\alpha^\star$,
\[
\alpha^\star\ \ge\ \alpha_2+\frac{(\alpha_2-\alpha_1)\,\bar D_{\alpha_2}}{\bar D_{\alpha_1}-\bar D_{\alpha_2}} ;
\]
\item $\epsilon\mapsto \bar D_{\alpha^\star-\epsilon}/\epsilon$ is nondecreasing on $(0,\alpha^\star]$, so it has a limit
$\mathfrak s\ge0$ as $\epsilon\downarrow0$, and $D_{\alpha^\star-\epsilon}(P,Q)\le D_0(P,Q)\,\epsilon/\{\alpha^\star(1-\alpha^\star)\}$
for all $\epsilon\in[0,\alpha^\star]$. Assumption~\ref{ass:margin} holds with $\gamma=1$ if and only if $\mathfrak s>0$, and then
it holds with $c_0=\mathfrak s$.
\end{enumerate}
\end{lemma}

\begin{proof}[Proof of Lemma~\ref{lem:pathconvex}]
A mixture of two trimmings at two levels is a trimming at the level in between, and the norm is convex; everything
else follows from convexity and the location of the zeros of the path.

We show first that $\bar D$ is convex. Let $\alpha_0,\alpha_1\in[0,1)$, $\lambda\in[0,1]$ and $\alpha_\lambda=(1-\lambda)\alpha_0+\lambda\alpha_1$.
If $h_0$ and $h_1$ are admissible weights at the levels $\alpha_0$ and $\alpha_1$, then $(1-\lambda)h_0+\lambda h_1$
lies between $0$ and $1$ and has mean $1-\alpha_\lambda$ under $P$, so
$(1-\lambda)\bar{\mathcal C}_{\alpha_0}(P)+\lambda \bar{\mathcal C}_{\alpha_1}(P)\subset \bar{\mathcal C}_{\alpha_\lambda}(P)$, and likewise for $Q$. For
$a_i\in \bar{\mathcal C}_{\alpha_i}(P)$ and $b_i\in \bar{\mathcal C}_{\alpha_i}(Q)$, the points $(1-\lambda)a_0+\lambda a_1$ and
$(1-\lambda)b_0+\lambda b_1$ therefore lie in the sets at level $\alpha_\lambda$, and the triangle inequality gives
\[
\bar D_{\alpha_\lambda}\le\big\|(1-\lambda)(a_0-b_0)+\lambda(a_1-b_1)\big\|_\Hh\le(1-\lambda)\|a_0-b_0\|_\Hh+\lambda\|a_1-b_1\|_\Hh .
\]
Taking the infimum over $a_0,a_1,b_0,b_1$ gives $\bar D_{\alpha_\lambda}\le(1-\lambda)\bar D_{\alpha_0}+\lambda \bar D_{\alpha_1}$. The path
$D_\alpha$ is nonnegative and nonincreasing, and so is $1-\alpha$, so their product $\bar D$ is nonincreasing.

From here on $0<\alpha^\star<1$, so that $\bar D_{\alpha^\star}=0$ and $\bar D>0$ on $[0,\alpha^\star)$ by
\eqref{eq:zeroeq}. The path is strictly
decreasing on $[0,\alpha^\star]$: for $\alpha_1<\alpha_2\le\alpha^\star$ write
$\alpha_2=(1-\lambda)\alpha_1+\lambda\alpha^\star$ with $\lambda=(\alpha_2-\alpha_1)/(\alpha^\star-\alpha_1)\in(0,1]$;
convexity gives $\bar D_{\alpha_2}\le(1-\lambda)\bar D_{\alpha_1}<\bar D_{\alpha_1}$, since $\bar D_{\alpha_1}>0$.

For the secant bound of (i), let $\alpha_1<\alpha_2<\alpha^\star$, so that $\bar D_{\alpha_1}>\bar D_{\alpha_2}>0$, and let
$\alpha\in(\alpha_2,1)$. Since $\alpha_2$ lies between $\alpha_1$ and $\alpha$, convexity gives
$(\alpha-\alpha_1)\bar D_{\alpha_2}\le(\alpha-\alpha_2)\bar D_{\alpha_1}+(\alpha_2-\alpha_1)\bar D_{\alpha}$, that is,
\[
\bar D_{\alpha}\ \ge\ \bar D_{\alpha_2}-(\alpha-\alpha_2)\,\frac{\bar D_{\alpha_1}-\bar D_{\alpha_2}}{\alpha_2-\alpha_1} .
\]
The right side, the secant through the two points, is positive for $\alpha<\alpha_s=\alpha_2+(\alpha_2-\alpha_1)\bar D_{\alpha_2}
/\{\bar D_{\alpha_1}-\bar D_{\alpha_2}\}$. So $\bar D>0$ on $[\alpha_2,\alpha_s)\cap[0,1)$, and since $\bar D_{\alpha^\star}=0$ with
$\alpha^\star<1$, this forces $\alpha^\star\ge\alpha_s$.

The slope at $\alpha^\star$ in (ii) comes from the same mixture argument with one endpoint at $\alpha^\star$. Put $r(\epsilon)=\bar D_{\alpha^\star-\epsilon}/\epsilon$ for
$\epsilon\in(0,\alpha^\star]$. For $0<\epsilon_1<\epsilon_2\le\alpha^\star$ the level $\alpha^\star-\epsilon_1$ is the
mixture of $\alpha^\star-\epsilon_2$ and $\alpha^\star$ with weights $\epsilon_1/\epsilon_2$ and $1-\epsilon_1/\epsilon_2$,
so convexity and $\bar D_{\alpha^\star}=0$ give $\bar D_{\alpha^\star-\epsilon_1}\le(\epsilon_1/\epsilon_2)\,\bar D_{\alpha^\star-\epsilon_2}$,
that is, $r(\epsilon_1)\le r(\epsilon_2)$. Hence $r$ is nondecreasing and nonnegative, its limit $\mathfrak s\ge0$ at $0$ exists, and
$\mathfrak s\le r(\epsilon)\le r(\alpha^\star)=D_0(P,Q)/\alpha^\star$ for every $\epsilon$. The upper bound gives
$D_{\alpha^\star-\epsilon}=\bar D_{\alpha^\star-\epsilon}/(1-\alpha^\star+\epsilon)\le D_0\,\epsilon/\{\alpha^\star(1-\alpha^\star)\}$
for $\epsilon\in(0,\alpha^\star]$, and at $\epsilon=0$ both sides vanish.

For the equivalence in (ii), if $\mathfrak s>0$, then $\bar D_{\alpha^\star-\epsilon}\ge \mathfrak s\epsilon$, and since $1-\alpha^\star+\epsilon\le1$,
$D_{\alpha^\star-\epsilon}\ge \bar D_{\alpha^\star-\epsilon}\ge \mathfrak s\epsilon$ for $\epsilon\in(0,\alpha^\star]$; at $\epsilon=0$ the
bound is $0\ge0$. So Assumption~\ref{ass:margin} holds with $\gamma=1$ and $c_0=\mathfrak s$. Conversely, if it holds with
$\gamma=1$ and some $c_0>0$, then $r(\epsilon)=(1-\alpha^\star+\epsilon)D_{\alpha^\star-\epsilon}/\epsilon\ge(1-\alpha^\star)c_0$
for every $\epsilon$, so $\mathfrak s\ge(1-\alpha^\star)c_0>0$. The bound $c_0\epsilon^\gamma\le D_0\,\epsilon/\{\alpha^\star(1-\alpha^\star)\}$
for all small $\epsilon>0$ rules out every exponent $\gamma<1$, which is why Assumption~\ref{ass:margin} takes
$\gamma\ge1$.
\end{proof}

The convexity and the monotonicity of $\bar D$, and $\bar D_{\alpha^\star}=0$, need only (A1) and (A2), since a
common trimming has one embedding; so does part (ii), which uses nothing else, while part (i) also needs
$\bar D>0$ below $\alpha^\star$, that is (A3).

\subsection{The rate of the lower bound}
\label{sec:lorate}

\begin{theorem}[Rate of the lower bound under the margin condition]
\label{thm:lorate}
Let Assumptions~\ref{ass:regularity} and~\ref{ass:margin} hold, let $\alpha^\star<1$, and compute
$\widehat\alpha$ as in Proposition~\ref{prop:loconsistency}: on a grid $\alpha_1<\dots<\alpha_J$ fixed in
advance, with a witness fitted at each level by a minimizer of \eqref{eq:qp} on the batches before the one bet on, cut
points satisfying \eqref{eq:cutrank}, and the fractions \eqref{eq:fractions} with $\lambda_{\max}=\tfrac12$,
$v_0^2\le4$ and the running variance floored at a constant $v^2_{\min}\in(0,4]$. Let $B$ be the number of
batches bet on, $N$ the number of pairs and $r_1=1/n_{<1}+1/m_{<1}$, where $n_{<1},m_{<1}$ are the sizes of batch $0$,
and suppose $r_1\le1$. Put $\beta_\star=(1-\alpha^\star)^{-1}$, $L_\delta=\log(1/\delta)$ and, for $\eta\in(0,1)$ and
$J'\ge1$,
\begin{gather}
T_{\mathrm{fit}}(\eta)=8\Big[2(1+\sqrt2)\kappa+\Delta_\phi\sqrt{\tfrac12\log(4B/\eta)}\,\Big]\sqrt{r_1},\notag\\
T_{\mathrm{bet}}(\eta,J')=2\Delta_\phi\sqrt{\frac2N}\,\Big[\Big(1+\frac{16}{v^2_{\min}}\Big)\sqrt{L_\delta}
+\frac2{v_{\min}}\sqrt{\log(2J'/\eta)}\,\Big].
\label{eq:loconst}
\end{gather}
For $\epsilon\in(0,\alpha^\star]$ let $a_\epsilon=\max(\{0\}\cup\{\alpha_j:\alpha_j\le\alpha^\star-\epsilon\})$, let
$\Delta_{\min}$ be the smallest gap between consecutive grid levels ($\Delta_{\min}=\infty$ if $J=1$), and let
$J_\epsilon=\min\{J,2(1+\epsilon/\Delta_{\min})\}$.
\begin{enumerate}[label=(\roman*),leftmargin=*]
\item If $c_0\epsilon^\gamma\ge\beta_\star\max\{T_{\mathrm{fit}}(\eta),T_{\mathrm{bet}}(\eta,J_\epsilon)\}$, then
$\Pp(\widehat\alpha<a_\epsilon)\le\eta$.
\item Let $t_{n,m}=\beta_\star\max\{T_{\mathrm{fit}}(\eta),T_{\mathrm{bet}}(\eta,J)\}$ and let
$\Delta_{\mathcal G}=\sup_{a\in[0,\alpha^\star]}\big(a-\max\{g\in\{0,\alpha_1,\dots,\alpha_J\}:g\le a\}\big)$ be the mesh of the grid
below $\alpha^\star$. Then, with probability at least $1-\delta-\eta$,
\[
\alpha^\star-\Delta_{\mathcal G}-\Big(\frac{t_{n,m}}{c_0}\Big)^{1/\gamma}\ \le\ \widehat\alpha\ \le\ \alpha^\star .
\]
\item With the batches of Section~\ref{sec:betting}, their boundaries rounded to the nearest integers, and
$n\wedge m\ge10$, one has $r_1\le\min\{1,12/(n\wedge m)\}$ and $N\ge(n\wedge m)/2$, so that, with an absolute
constant in the $O$,
\[
\Big(\frac{t_{n,m}}{c_0}\Big)^{1/\gamma}
=O\bigg(\bigg(\frac{\kappa}{c_0(1-\alpha^\star)v^2_{\min}}
\sqrt{\frac{1+\log(1/\delta)+\log(J/\eta)}{n\wedge m}}\bigg)^{1/\gamma}\bigg).
\]
\end{enumerate}
\end{theorem}

The error is of order $(n\wedge m)^{-1/(2\gamma)}$, plus the mesh, with the factor $1/v^2_{\min}$ of the betting
constants and a factor $\log J$, which part (i) removes when the grid is not finer than the rate: if $\Delta_{\min}\ge\epsilon'_{n,m}=(\beta_\star\max\{T_{\mathrm{fit}}(\eta),T_{\mathrm{bet}}(\eta,4)\}/c_0)^{1/\gamma}$,
then $J_{\epsilon'_{n,m}}\le4$, and with probability at least $1-\delta-\eta$,
$\alpha^\star-\Delta_{\mathcal G}-\epsilon'_{n,m}\le\widehat\alpha\le\alpha^\star$, with no $\log J$. The proof needs no
density of the witness at the cut points, unlike Theorem~\ref{cor:trimrate}, which uses \eqref{eq:qdens}, and it does
not use (A3), whose role the margin takes over.

\begin{proof}[Proof of Theorem~\ref{thm:lorate}]
We run the lower side of the proof of Proposition~\ref{prop:loconsistency} at the level $\alpha^\star-\epsilon$, with the
margin in place of a fixed positive distance, and keep the constants. One event, which does not depend on the level,
makes every fitted ramp good, and only the bets need a union over the levels.

Fix $\epsilon$ as in (i), and put $D_*=c_0\epsilon^\gamma$, $\tau_*=D_*/(8\beta_\star)$ and
$\varrho_*=D_*/(2\beta_\star\Delta_\phi)$: $D_*$ is the margin at distance $\epsilon$ from $\alpha^\star$, $\tau_*$ the
fluctuation allowed to each fitting batch, and $\varrho_*$ the excess that the bets will see. The two thresholds of
(i) are these choices solved for $D_*$: $T_{\mathrm{fit}}$ makes the fitting suprema smaller than $\tau_*$, and
$T_{\mathrm{bet}}$ makes an excess $\varrho_*$ win the bets, each with probability $\eta/2$ to spare. Every grid level $\alpha_j\le\alpha^\star-\epsilon$ has
$\beta_j=(1-\alpha_j)^{-1}\le\beta_\star$, and with $x_j=(\alpha^\star-\alpha_j)/\epsilon\ge1$,
Assumption~\ref{ass:margin} at $\epsilon x_j\in[\epsilon,\alpha^\star]$ gives
$D_{\alpha_j}\ge c_0\epsilon^\gamma x_j^\gamma=D_*x_j^\gamma\ge D_*$. Moreover $D_*\le D_{\alpha^\star-\epsilon}\le
D_0=\MMD(P,Q)=\|\E[\phi(X)-\phi(Y)]\|_\Hh\le\Delta_\phi$, as the path is nonincreasing, so $\Delta_\phi>0$. By
\eqref{eq:downclosed}, if $H_{0,\alpha_j}$ is rejected at every grid level $\alpha_j\le\alpha^\star-\epsilon$, then
$\widehat\alpha\ge a_\epsilon$, which also holds trivially when no grid level is that small. Hence
\[
\{\widehat\alpha<a_\epsilon\}\subseteq\bigcup_{j:\,\alpha_j\le\alpha^\star-\epsilon}
\{H_{0,\alpha_j}\text{ is not rejected}\}.
\]

The fits are controlled by the suprema $\Gamma_b,\Gamma_b^-$ of the proof of Proposition~\ref{prop:loconsistency},
those of the proof of Theorem~\ref{thm:twosided} computed on batches $0,\dots,b-1$, which do not depend on $\alpha$. Let
$\mathcal E_{\mathrm{fit}}=\bigcap_{b=1}^B\{\Gamma_b\le\tau_*,\ \Gamma_b^-\le\tau_*\}$. We show that on
$\mathcal E_{\mathrm{fit}}$ every batch has a ramp at every level $\alpha_j\le\alpha^\star-\epsilon$, with a large
excess. Fix such a level $\alpha=\alpha_j$ and a batch $b$, and write $\beta=\beta_j$. The hypothesis of (i) gives
$D_*\ge\beta_\star T_{\mathrm{fit}}(\eta)\ge16(1+\sqrt2)\beta_\star\kappa\sqrt{r_1}$, and since $\Delta_\phi\le2\kappa$
and $r_b\le r_1\le1$,
\[
\beta\Delta_\phi r_b\le2\beta_\star\kappa\sqrt{r_1}\le\frac{D_*}8,
\qquad
\beta^2\kappa^2r_b\le\beta_\star^2\kappa^2r_1\le\frac{D_*^2}{256(1+\sqrt2)^2}\le\frac{7D_*^2}{64};
\]
on $\mathcal E_{\mathrm{fit}}$ moreover $\beta(\Gamma_b+\Gamma_b^-)\le2\beta_\star\tau_*=D_*/4$ and
$D_\alpha-\beta\Gamma_b^-\ge D_*-D_*/8=7D_*/8>0$. The last inequality is the hypothesis
$\beta\Gamma_b^-<D_\alpha$ under which the proof of Proposition~\ref{prop:loconsistency} establishes
\eqref{eq:fittedramp}, with no condition on the law of the witness near the cut points:
\[
\Delta_\phi\ge\widehat t_b-\widehat u_b\ge\theta(\widehat f;\widehat t_b,\widehat u_b)\ge\Theta_b
=D_\alpha-\beta(\Gamma_b+\Gamma_b^-)-\beta\Delta_\phi r_b-\frac{\beta^2\kappa^2r_b}{D_\alpha-\beta\Gamma_b^-},
\]
and the four bounds above give $\Theta_b\ge D_\alpha-\frac{D_*}4-\frac{D_*}8-\frac{7D_*^2/64}{7D_*/8}=D_\alpha-\frac{D_*}2
\ge\frac{D_\alpha}2>0$. So $\widehat t_b>\widehat u_b$, batch $b$ has a ramp $\rp_b$, and Lemma~\ref{lem:ramp}(ii),
applied under $P$ and $Q$ given batches $0,\dots,b-1$ to the direction $\widehat f$ and the cut points $\widehat
t_b,\widehat u_b$, gives $\beta(\widehat t_b-\widehat u_b)(\E_P\rp_b-\E_Q\rp_b-\alpha)\ge\theta(\widehat f;\widehat
t_b,\widehat u_b)\ge\Theta_b$, hence
\begin{equation}
\E_P\rp_b-\E_Q\rp_b-\alpha_j\ \ge\ \frac{\Theta_b}{\beta\Delta_\phi}\ \ge\ \frac{D_{\alpha_j}}{2\beta_\star\Delta_\phi}
\ \eqqcolon\ \varrho_j\ \ge\ \varrho_*x_j^\gamma .
\label{eq:excessj}
\end{equation}
The fitting event is likely. By Lemma~\ref{lem:hinge} with $\varphi^1_x=\varphi^2_y=(\cdot)_+$, $r=1$ and
$\rho_1=\rho_2=\kappa$, as in the proof of Theorem~\ref{thm:twosided}, $\E\Gamma_b$ and $\E\Gamma_b^-$ are at most
$2\kappa\sqrt{r_b}+2\kappa(e_{n_{<b}}+e_{m_{<b}})\le2(1+\sqrt2)\kappa\sqrt{r_b}\le\mu_1=2(1+\sqrt2)\kappa\sqrt{r_1}$, by
$e_k\le k^{-1/2}$ and $n^{-1/2}+m^{-1/2}\le\{2(1/n+1/m)\}^{1/2}$, and, the bounded differences being $\Delta_\phi$, each
exceeds its mean by more than $y$ with probability at most $e^{-2y^2/(\Delta_\phi^2r_b)}\le e^{-2y^2/(\Delta_\phi^2r_1)}$.
Divided by $8\beta_\star$, the hypothesis $D_*\ge\beta_\star T_{\mathrm{fit}}(\eta)$ reads
$\tau_*-\mu_1\ge\Delta_\phi\{r_1\log(4B/\eta)/2\}^{1/2}$, so each of the $2B$ suprema exceeds $\tau_*$ with probability at
most $\eta/(4B)$, and $\Pp(\mathcal E_{\mathrm{fit}}^c)\le\eta/2$.

We turn to the bets, where only a union over the levels is needed. Fix a grid level $\alpha_j\le\alpha^\star-\epsilon$. On $\mathcal E_{\mathrm{fit}}$ every batch has a
ramp, and by \eqref{eq:excessj} every pair has conditional mean $\Delta_k\ge\alpha_j+\varrho_j$ in the notation of the
proof of Theorem~\ref{thm:betting}. The sample is large enough for \eqref{eq:betbounds}: as $v^2_{\min}\le4$,
$\Delta_\phi\ge D_*\ge\beta_\star T_{\mathrm{bet}}(\eta,J_\epsilon)\ge10\Delta_\phi(2L_\delta/N)^{1/2}$, so $N\ge200L_\delta\ge
2L_\delta/\lambda_{\max}^2$. The first bound of \eqref{eq:betbounds}, with $\mathcal E=\mathcal E_{\mathrm{fit}}$ and
$\varrho=\varrho_j$, gives
\[
\Pp\big(\mathcal E_{\mathrm{fit}}\cap\{H_{0,\alpha_j}\text{ is not rejected}\}\big)\le\varepsilon_N(\varrho_j)
=\exp\Big\{-\frac{v^2_{\min}}{4L_\delta}\big(\varrho_jS-C\big)_+^2\Big\},
\]
with $S=(NL_\delta/2)^{1/2}$ and $C=L_\delta(1+16/v^2_{\min})$. The hypothesis $D_*\ge\beta_\star T_{\mathrm{bet}}(\eta,J_\epsilon)$, divided by $2\beta_\star\Delta_\phi$ and multiplied
by $S$, reads $\varrho_*S\ge C+2v_{\min}^{-1}\{L_\delta\log(2J_\epsilon/\eta)\}^{1/2}$, so
$a=\varrho_*S-C\ge2v_{\min}^{-1}\{L_\delta\log(2J_\epsilon/\eta)\}^{1/2}>0$. Since $x_j\ge1$, $\gamma\ge1$ and
$\varrho_*S\ge a$, \eqref{eq:excessj} gives $\varrho_jS-C\ge\varrho_*x_jS-C=a+\varrho_*S(x_j-1)\ge ax_j$, hence
$\varepsilon_N(\varrho_j)\le e^{-Ax_j^2}$ with $A=v^2_{\min}a^2/(4L_\delta)\ge\log(2J_\epsilon/\eta)>\log2$.

It remains to add over the levels, and the polynomial growth of the excess away from $\alpha^\star$ keeps the sum
within a factor $J_\epsilon$ of $e^{-A}$. If $J_\epsilon=J$, the at most $J$ terms $e^{-Ax_j^2}\le e^{-A}$ add up to at most
$Je^{-A}\le\eta/2$. Otherwise $J_\epsilon=2(1+\epsilon/\Delta_{\min})$; the $x_j$ are at least $\Delta_{\min}/\epsilon$ apart, so at most
$1+\epsilon/\Delta_{\min}$ of them lie in each interval $[1+k,2+k)$, $k=0,1,\dots$, and as $(1+k)^2\ge1+3k$ and
$e^{-3A}\le\tfrac12$,
\[
\sum_{j:\,\alpha_j\le\alpha^\star-\epsilon}e^{-Ax_j^2}\le\Big(1+\frac\epsilon{\Delta_{\min}}\Big)\sum_{k\ge0}e^{-A(1+3k)}
\le2\Big(1+\frac\epsilon{\Delta_{\min}}\Big)e^{-A}=J_\epsilon e^{-A}\le\frac\eta2 .
\]
With the inclusion above, $\Pp(\widehat\alpha<a_\epsilon)\le\Pp(\mathcal E_{\mathrm{fit}}^c)
+\sum_j\Pp(\mathcal E_{\mathrm{fit}}\cap\{H_{0,\alpha_j}\text{ is not rejected}\})\le\eta$, which is (i).

For (ii), put $\epsilon_{n,m}=(t_{n,m}/c_0)^{1/\gamma}$. If $\epsilon_{n,m}>\alpha^\star$ the lower inequality holds
because $\widehat\alpha\ge0$. Otherwise $\epsilon=\epsilon_{n,m}$ satisfies the hypothesis of (i), because
$J_\epsilon\le J$ and $T_{\mathrm{bet}}$ increases in $J'$, so with probability at least $1-\eta$,
$\widehat\alpha\ge a_{\epsilon_{n,m}}\ge\alpha^\star-\epsilon_{n,m}-\Delta_{\mathcal G}$, the last step by the definition of
$\Delta_{\mathcal G}$. Corollary~\ref{cor:lower} gives $\widehat\alpha\le\alpha^\star$ with probability at least $1-\delta$,
and a union bound gives (ii).

For (iii), boundaries at the nearest integers to $n/4$ and $n/2$ give batch $0$ at least $(n-2)/4$ observations $X$,
batch $1$ at least $n/4-1$ and batch $2$ at least $n/2-\tfrac12$, and likewise for $Y$. Hence $r_1\le8/(n\wedge m-2)$,
which is at most $1$ for $n\wedge m\ge10$ and at most $12/(n\wedge m)$ for $n\wedge m\ge6$, and
$N\ge\tfrac34(n\wedge m)-\tfrac32\ge(n\wedge m)/2$ for $n\wedge m\ge6$. Inserting these bounds, $B=2$,
$\Delta_\phi\le2\kappa$, and $1+16/v^2_{\min}\le20/v^2_{\min}$ and $2/v_{\min}\le4/v^2_{\min}$, which
hold as $v^2_{\min}\le4$, into \eqref{eq:loconst} bounds $t_{n,m}$ by an absolute constant times
$\beta_\star\kappa v_{\min}^{-2}\{(1+L_\delta+\log(J/\eta))/(n\wedge m)\}^{1/2}$.
\end{proof}

\subsection{Touching contaminations}
\label{sec:newproofs}

This subsection and the six that follow prove Propositions~\ref{prop:newslope}--\ref{prop:newgauss} of the main
text, in a more general form, state and prove the rates under a general margin (Proposition~\ref{prop:newrates}), and
state and prove the converse Proposition~\ref{prop:newconverse}. Throughout, $P\wedge Q$ has mass $1-\alpha^\star$,
$P_0=(P-P\wedge Q)/\alpha^\star$ and $Q_0=(Q-P\wedge Q)/\alpha^\star$, so that $P=P\wedge Q+\alpha^\star P_0$,
$Q=P\wedge Q+\alpha^\star Q_0$, $P_0$ is concentrated on $A=\{p>q\}$, $Q_0$ on its complement, and
$|P-Q|=\alpha^\star(P_0+Q_0)$. The rescaled path $\bar D$ and the sets $\bar{\mathcal C}_\alpha$ are those of
Section~\ref{sec:generalmargin}, and for $\alpha\le\alpha^\star$, $\bar D_{\alpha}\le D_\alpha(P,Q)\le
\bar D_{\alpha}/(1-\alpha^\star)$.

On $\R^d$ we use translation-invariant kernels $k(x,y)=k_0(x-y)$ with $k_0$ continuous, integrable and positive
definite, and the Fourier transform $\widehat g(\xi)=\int g(x)e^{-i\xi\cdot x}\,dx$. Then $\widehat{k_0}\ge0$ is
integrable, $k_0(z)=(2\pi)^{-d}\int\widehat{k_0}(\xi)e^{i\xi\cdot z}\,d\xi$, and by
\citet[Theorem~10.12]{wendland2005scattered} the RKHS consists of the continuous $f\in L^2(\R^d)$ with
$\int|\widehat f|^2/\widehat{k_0}<\infty$, with
\begin{equation}
\|f\|_\Hh^2=(2\pi)^{-d}\int\frac{|\widehat f(\xi)|^2}{\widehat{k_0}(\xi)}\,d\xi .
\label{eq:newnative}
\end{equation}
For a finite signed measure $\nu$ with Fourier transform $\widehat\nu(\xi)=\int e^{-i\xi\cdot x}\,d\nu(x)$, Fubini's
theorem and the inversion formula give
\begin{equation}
\Big\|\int\phi\,d\nu\Big\|_\Hh^2=\iint k_0(x-y)\,d\nu(x)\,d\nu(y)=(2\pi)^{-d}\int|\widehat\nu(\xi)|^2\,\widehat{k_0}(\xi)\,d\xi .
\label{eq:newbochner}
\end{equation}

The first result needs no topology on $\mathcal X$. By (A1), (A2) and Pettis' theorem, $\phi$ is Borel measurable
into the separable space $\Hh$, so its image laws $\phi_\#P_0$ and $\phi_\#Q_0$ are Borel probability measures on a
separable metric space, whose supports are closed, nonempty and of full mass.

\begin{proposition}[Slope bound]
\label{prop:newslopegeneral}
Let (A1) and (A2) hold and let $0<\alpha^\star<1$. Then
$\lim_{\epsilon\downarrow0}\bar D_{\alpha^\star-\epsilon}/\epsilon\le\dist\big(\supp\phi_\#P_0,\supp\phi_\#Q_0\big)$.
\end{proposition}

\begin{proof}
We trim at the level $\alpha^\star-\epsilon$ by adding to the common part a little of each contamination near two
chosen feature vectors. Fix $z\in\supp\phi_\#P_0$, $z'\in\supp\phi_\#Q_0$ and $\varrho>0$, and let
$U=\phi^{-1}(B(z,\varrho))$ and $V=\phi^{-1}(B(z',\varrho))$, so that $m_U=P_0(U)>0$ and $m_V=Q_0(V)>0$. Write
$P_0^U=P_0(\cdot\cap U)/m_U$ and $Q_0^V=Q_0(\cdot\cap V)/m_V$. For $0<\epsilon\le\alpha^\star\min(m_U,m_V)$ the measure
$P_\epsilon=P\wedge Q+\epsilon P_0^U$ satisfies $P_\epsilon\le P\wedge Q+\alpha^\star P_0=P$, because $\epsilon/m_U\le\alpha^\star$, and has
mass $1-(\alpha^\star-\epsilon)$, so $\int\phi\,dP_\epsilon\in \bar{\mathcal C}_{\alpha^\star-\epsilon}(P)$. Likewise,
$Q_\epsilon=P\wedge Q+\epsilon Q_0^V$ gives a point of $\bar{\mathcal C}_{\alpha^\star-\epsilon}(Q)$. The common part cancels, and
\begin{align*}
\bar D_{\alpha^\star-\epsilon}&\le\Big\|\int\phi\,dP_\epsilon-\int\phi\,dQ_\epsilon\Big\|_\Hh
=\epsilon\,\Big\|\iint\{\phi(x)-\phi(y)\}\,dP_0^U(x)\,dQ_0^V(y)\Big\|_\Hh\\
&\le\epsilon\,(\|z-z'\|_\Hh+2\varrho),
\end{align*}
since $\|\phi(x)-z\|_\Hh<\varrho$ on $U$ and $\|\phi(y)-z'\|_\Hh<\varrho$ on $V$. Dividing by $\epsilon$ and letting
$\epsilon\downarrow0$, and then $\varrho\downarrow0$, bounds the limit, which exists by Lemma~\ref{lem:pathconvex}(ii) (that part
of its proof uses only the convexity of $\bar D$ and $\bar D_{\alpha^\star}=0$, which hold under (A1) and (A2)), by $\|z-z'\|_\Hh$, and the infimum over $z$ and $z'$ gives the claim.
\end{proof}

\begin{proof}[Proof of Proposition~\ref{prop:newslope}]
The quantity bounded in the proposition is $\lim \bar D_{\alpha^\star-\epsilon}/\epsilon$. If $\phi$ is continuous and
$x\in\supp P_0$, then for every $\varrho>0$ the set $\phi^{-1}(B(\phi(x),\varrho))$ is an open neighborhood of $x$ and has
positive $P_0$-mass, so $\phi(x)\in\supp\phi_\#P_0$. Likewise $\phi(y)\in\supp\phi_\#Q_0$ for $y\in\supp Q_0$, and
Proposition~\ref{prop:newslopegeneral} gives the bound. If the supports meet, take $y=x$: the limit is $0$, and by
Lemma~\ref{lem:pathconvex}(ii) Assumption~\ref{ass:margin} does not hold with $\gamma=1$.

For the last claim let $\mathcal X$ be an open connected subset of $\R^d$ and let $p,q$ be continuous. The density
$(p-q)_+/\alpha^\star$ of $P_0$ is continuous, so $\supp P_0$ is the closure $F_P$ in $\mathcal X$ of the open set
$\{p>q\}$: a neighborhood of a point of $F_P$ meets $\{p>q\}$ in a nonempty open set, which has positive Lebesgue
measure and on which the density is positive, while a point outside $F_P$ has a neighborhood of $P_0$-mass $0$.
Likewise $\supp Q_0$ is the closure $F_Q$ of $\{q>p\}$. Both sets are nonempty, as $\int(p-q)_+=\int(q-p)_+=\alpha^\star>0$.
If $F_P\cap F_Q$ were empty, the open set $\mathcal X\setminus(F_P\cup F_Q)$ would lie in $\{p=q\}$, because
$\{p\neq q\}\subseteq F_P\cup F_Q$, and so would be empty. Then $\mathcal X$ would be the union of the disjoint nonempty
sets $F_P$ and $F_Q$, closed in $\mathcal X$, against its connectedness. For two distinct finite Gaussian mixtures,
$p-q$ is real-analytic on $\R^d$ and not identically zero, so its zero set has empty interior, and $\alpha^\star<1$
because both densities are positive.
\end{proof}

When $\operatorname{supp}P_0$ is compact, the bound is an equivalence for continuous kernels whose RKHS contains the
smooth functions with compact support.

\begin{proposition}[Disjoint contaminations give a linear margin]
\label{prop:newgammaone}
Let $\mathcal X=\R^d$ and $k(x,y)=k_0(x-y)$ with $\widehat{k_0}(\xi)\ge c(1+|\xi|^2)^{-s}$ for some $s>d/2$ and $c>0$,
as for the Mat\'ern kernels, let $0<\alpha^\star<1$ and let $\supp P_0$ be compact. If $\supp P_0\cap\supp Q_0=\varnothing$,
there is $f\in\Hh$ with $0\le f\le1$, $f=1$ on $\supp P_0$ and $f=0$ on $\supp Q_0$, and every such $f$ gives
$D_{\alpha^\star-\epsilon}(P,Q)\ge \bar D_{\alpha^\star-\epsilon}\ge\epsilon/\|f\|_\Hh$ for $\epsilon\in[0,\alpha^\star]$, that is,
Assumption~\ref{ass:margin} with $\gamma=1$ and $c_0=1/\|f\|_\Hh$. Together with Proposition~\ref{prop:newslope}, the
rescaled path then has a positive slope at $\alpha^\star$ if and only if $\supp P_0\cap\supp Q_0=\varnothing$.
\end{proposition}

\begin{proof}
A compact set and a disjoint closed set are at a positive distance $d_0$. Let $U$ be the open
$d_0/2$-neighborhood of $\supp P_0$, which is bounded, and $f=\mathbf 1_U*b_{d_0/3}$ with $b\in C_c^\infty$ a
probability density supported in the closed unit ball and $b_w=w^{-d}b(\cdot/w)$. Then $f$ is smooth with compact
support and $0\le f\le1$. Moreover $f=1$ on $\supp P_0$, since the closed ball of radius $d_0/3$ about such a point lies in $U$,
and $f=0$ on $\supp Q_0$, since the closed ball of radius $d_0/3$ about such a point stays at distance at least
$2d_0/3$ from $\supp P_0$ and so misses $U$. By \eqref{eq:newnative},
$\|f\|_\Hh^2\le c^{-1}(2\pi)^{-d}\int|\widehat f|^2(1+|\xi|^2)^s<\infty$, so $f\in\Hh$. For any such $f$,
$\E_Pf-\E_Qf=\alpha^\star(\E_{P_0}f-\E_{Q_0}f)=\alpha^\star$, so $\mathrm{err}(f)=0$, and Proposition~\ref{prop:newapproxgen}(ii)
below gives the bound. The equivalence follows from Lemma~\ref{lem:pathconvex}(ii) and
Proposition~\ref{prop:newslope}, as $\phi$ is continuous.
\end{proof}

This sharpens Example~\ref{ex:separated}, which also separates the core.

\subsection{Margins from approximate witnesses}

\begin{lemma}[Trimmed means of a bounded function]
\label{lem:newtrimmean}
Let $\alpha\in[0,1)$, let $P$ be a probability measure and $f$ a bounded measurable function. Every
$P'\in R_\alpha(P)$ satisfies
\[
\frac{\E_Pf-\alpha\sup f}{1-\alpha}\ \le\ \E_{P'}f\ \le\ \frac{\E_Pf-\alpha\inf f}{1-\alpha} .
\]
\end{lemma}

\begin{proof}
Write $P'=gP$ with $0\le g\le\beta$ and $\E_Pg=1$. Then $\E_{P'}f=\beta\E_Pf-\E_P[(\beta-g)f]$, where $\beta-g\ge0$ and
$\E_P(\beta-g)=\beta-1=\alpha\beta$, so $\alpha\beta\inf f\le\E_P[(\beta-g)f]\le\alpha\beta\sup f$.
\end{proof}

These are the crude forms of Lemma~\ref{lem:cvar}: a trimming can at worst remove the mass $\alpha$ at the largest, or
the smallest, value of $f$. We write $\operatorname{osc}(f)=\sup f-\inf f$, $\mathcal F=\{f\in\Hh:0\le f\le1\}$, and, for
$\epsilon\in[0,\alpha^\star]$,
\[
\omega_\Hh(\epsilon)=\sup\Big\{\frac{(\epsilon-\mathrm{err}(f))_+}{\|f\|_\Hh}:\ f\in\mathcal F,\ f\neq0\Big\},
\qquad
\mathcal K(\lambda)=\inf_{f\in\mathcal F}\big\{\mathrm{err}(f)+\lambda\|f\|_\Hh\big\} ,
\]
with $\sup\varnothing=0$.

\begin{proposition}[Approximation margin]
\label{prop:newapproxgen}
Let (A1) and (A2) hold.
\begin{enumerate}[label=(\roman*),leftmargin=*]
\item For every $f\in\Hh$, $f\neq0$, and $\alpha\in[0,1)$,
$(1-\alpha)\|f\|_\Hh D_\alpha(P,Q)\ge\E_Pf-\E_Qf-\alpha\operatorname{osc}(f)$.
\item Let $0<\alpha^\star<1$. For every $f\in\mathcal F$, $\E_Pf-\E_Qf=\alpha^\star-\mathrm{err}(f)$, and if $f\neq0$, then
$\bar D_{\alpha^\star-\epsilon}\ge\{\epsilon-\mathrm{err}(f)\}/\|f\|_\Hh$ for every $\epsilon\in[0,\alpha^\star]$.
\item $\omega_\Hh$ is nondecreasing, $D_{\alpha^\star-\epsilon}(P,Q)\ge \bar D_{\alpha^\star-\epsilon}\ge\omega_\Hh(\epsilon)\ge
\epsilon/\{2\mathcal W(\epsilon/2)\}$, read as $0$ when $\mathcal W(\epsilon/2)=\infty$, and for $\lambda\ge0$ and
$\epsilon\in(0,\alpha^\star]$, $\omega_\Hh(\epsilon)>\lambda$ if and only if $\mathcal K(\lambda)<\epsilon$.
\end{enumerate}
\end{proposition}

\begin{proof}
Part (i) bounds each trimmed mean by Lemma~\ref{lem:newtrimmean}, and part (ii) evaluates the bound for
$[0,1]$-valued $f$. For (i), $f$ is bounded, by $\kappa\|f\|_\Hh$. For $P'\in R_\alpha(P)$ and $Q'\in R_\alpha(Q)$,
Lemma~\ref{lem:newtrimmean} gives $\E_{P'}f-\E_{Q'}f\ge\{\E_Pf-\E_Qf-\alpha\operatorname{osc}(f)\}/(1-\alpha)$, and by the
reproducing property and the Cauchy--Schwarz inequality $\MMD(P',Q')\ge(\E_{P'}f-\E_{Q'}f)/\|f\|_\Hh$, and the infimum over
$P'$ and $Q'$ gives (i).

For (ii), $\int\mathbf 1_A\,d(P-Q)=\int(p-q)_+=\alpha^\star$, so $\E_Pf-\E_Qf=\alpha^\star-\int(\mathbf 1_A-f)(p-q)\,d(P+Q)$. On $A$
both $1-f$ and $p-q$ are nonnegative, and off $A$ both $-f$ and $p-q$ are nonpositive, so the integrand is
$|\mathbf 1_A-f|\,|p-q|$ and the integral is $\mathrm{err}(f)$. As $\operatorname{osc}(f)\le1$, part (i) at
$\alpha=\alpha^\star-\epsilon$ gives $\|f\|_\Hh \bar D_{\alpha^\star-\epsilon}\ge\alpha^\star-\mathrm{err}(f)-(\alpha^\star-\epsilon)$.

For (iii), each map $\epsilon\mapsto(\epsilon-\mathrm{err}(f))_+/\|f\|_\Hh$ is nondecreasing, and so is their supremum,
which is at most $\bar D_{\alpha^\star-\epsilon}$ by (ii), and $\bar D_{\alpha}\le D_\alpha$. If $\mathcal W(\epsilon/2)<\infty$, every
$f\in\mathcal F$ with $\mathrm{err}(f)\le\epsilon/2$ is nonzero, because $\mathrm{err}(0)=\alpha^\star>\epsilon/2$, and has
$(\epsilon-\mathrm{err}(f))/\|f\|_\Hh\ge\epsilon/(2\|f\|_\Hh)$. The supremum over these $f$ is
$\epsilon/\{2\mathcal W(\epsilon/2)\}$. If $\mathcal K(\lambda)<\epsilon$, some $f\in\mathcal F$ has
$\mathrm{err}(f)+\lambda\|f\|_\Hh<\epsilon$, and it is nonzero because $\mathrm{err}(0)=\alpha^\star\ge\epsilon$, so
$\omega_\Hh(\epsilon)\ge(\epsilon-\mathrm{err}(f))/\|f\|_\Hh>\lambda$. Conversely, $\omega_\Hh(\epsilon)>\lambda$ gives a
nonzero $f\in\mathcal F$ with $\epsilon-\mathrm{err}(f)>\lambda\|f\|_\Hh$, that is, $\mathcal K(\lambda)<\epsilon$.
\end{proof}

\begin{proof}[Proof of Proposition~\ref{prop:newapprox}]
This is Proposition~\ref{prop:newapproxgen}(ii) and (iii), with $\bar D_{\alpha^\star-\epsilon}=(1-\alpha^\star+\epsilon)D_{\alpha^\star-\epsilon}$.
\end{proof}

\subsection{Rates under a general margin}

The rate of Theorem~\ref{thm:loratemain} extends to any margin, polynomial or not. For a nondecreasing $\omega$ we
write $\omega^{\leftarrow}(t)=\sup\{\epsilon\in[0,\alpha^\star]:\omega(\epsilon)\le t\}$.

\begin{proposition}[Rates under a general margin]
\label{prop:newrates}
Let Assumption~\ref{ass:regularity} hold, let $\alpha^\star<1$, compute $\widehat\alpha$ as in
Theorem~\ref{thm:loratemain}, and let $\omega:[0,\alpha^\star]\to[0,\infty)$ be nondecreasing with
$(1-\alpha^\star+\epsilon)\,D_{\alpha^\star-\epsilon}(P,Q)\ge\omega(\epsilon)$ for every $\epsilon\in[0,\alpha^\star]$.
\begin{enumerate}[label=(\roman*),leftmargin=*]
\item Parts (i) and (ii) of Theorem~\ref{thm:lorate} hold with $c_0\epsilon^\gamma$ replaced by
$\omega(\epsilon)$ and $(t_{n,m}/c_0)^{1/\gamma}$ by $\omega^{\leftarrow}(t_{n,m})$.
\item If moreover $\omega(\epsilon)\ge\exp(-c'/\epsilon)$ for $0<\epsilon\le\epsilon_0$, for some $c'>0$ and
$\epsilon_0\in(0,\alpha^\star]$, and the grid has at most $(n\wedge m)^b$ levels for a fixed $b>0$, then for every
$\eta\in(0,1)$ and every $c>2c'$ there is $n_0$, depending on $c,c',\epsilon_0,b,\eta,\delta,\kappa,\alpha^\star$ and
$v^2_{\min}$, such that
$\Pp\{\alpha^\star-\Delta_{\mathcal G}-c/\log(n\wedge m)\le\widehat\alpha\le\alpha^\star\}\ge1-\delta-\eta$
whenever $n\wedge m\ge n_0$.
\item Under the hypotheses of Proposition~\ref{prop:newgauss}, with a grid as in (ii) and no $\omega$ assumed, the
conclusion of (ii) holds for every $c>\pi^2\sigma^2\ell_0/4$.
\end{enumerate}
\end{proposition}

\begin{proof}[Proof of Proposition~\ref{prop:newrates}(i)]
In the proof of Theorem~\ref{thm:lorate}, the margin enters only through $D_*=c_0\epsilon^\gamma$ and two facts: $D_*\le
D_{\alpha^\star-\epsilon}$, and $D_{\alpha_j}\ge D_*x_j$ for the grid levels $\alpha_j\le\alpha^\star-\epsilon$, with
$x_j=(\alpha^\star-\alpha_j)/\epsilon\ge1$. That proof writes $D_{\alpha_j}\ge D_*x_j^\gamma$ and then uses only
$x_j^\gamma\ge x_j$, in the step $\varrho_jS-C\ge\varrho_*x_jS-C$. Put $D_*=\omega(\epsilon)$. Then
$D_*\le \bar D_{\alpha^\star-\epsilon}\le D_{\alpha^\star-\epsilon}$. By Lemma~\ref{lem:pathconvex}(ii), $\epsilon'\mapsto
\bar D_{\alpha^\star-\epsilon'}/\epsilon'$ is nondecreasing on $(0,\alpha^\star]$, and $\epsilon x_j=\alpha^\star-\alpha_j\in[\epsilon,\alpha^\star]$,
so
\[
D_{\alpha_j}\ \ge\ \bar D_{\alpha_j}\ =\ \bar D_{\alpha^\star-\epsilon x_j}\ \ge\ x_j\,\bar D_{\alpha^\star-\epsilon}\ \ge\ x_jD_* .
\]
The rest of the proof of part (i) of Theorem~\ref{thm:lorate} uses $D_*\le\Delta_\phi$, which holds since
$D_{\alpha^\star-\epsilon}\le D_0\le\Delta_\phi$, and the hypothesis on $D_*$, and is unchanged. For part (ii), every
$\epsilon\in(\omega^{\leftarrow}(t_{n,m}),\alpha^\star]$ has $\omega(\epsilon)>t_{n,m}$ and $J_\epsilon\le J$, so part (i) gives
$\Pp(\widehat\alpha<\alpha^\star-\epsilon-\Delta_{\mathcal G})\le\Pp(\widehat\alpha<a_\epsilon)\le\eta$. Along a sequence
$\epsilon_k\downarrow\omega^{\leftarrow}(t_{n,m})$ these events increase to
$\{\widehat\alpha<\alpha^\star-\omega^{\leftarrow}(t_{n,m})-\Delta_{\mathcal G}\}$, whose probability is therefore at most $\eta$, and
Corollary~\ref{cor:lower} and a union bound finish the proof as in Theorem~\ref{thm:lorate}.
\end{proof}

Applied to a polynomial margin of $D$, part (i) loses a factor $(1-\alpha^\star)^{1/\gamma}$ against
Theorem~\ref{thm:lorate}, since a margin of $D$ gives one of $\bar D$ only after multiplication by $1-\alpha^\star$. Its use is
for margins that are not polynomial, such as the exponential margins of part (ii). Part (iii) is proved after
Theorem~\ref{thm:newgauss}.

\begin{proof}[Proof of Proposition~\ref{prop:newrates}(ii)]
Write $\underline n=n\wedge m$ and let $\underline n\ge10$. By Theorem~\ref{thm:lorate}(iii), the quantity $t_{n,m}$
of its part (ii) is at most an absolute constant times
$\beta_\star\kappa v_{\min}^{-2}\{(1+L_\delta+\log(J/\eta))/\underline n\}^{1/2}$, and with $J\le\underline n^b$
this is at most $C_t(\log\underline n/\underline n)^{1/2}$, with $C_t$ depending only on $\kappa$, $\alpha^\star$,
$v_{\min}$, $\delta$, $\eta$ and $b$. Let $0<t<\min\{1,\omega(\epsilon_0)\}$. As $\omega$ is nondecreasing, every
$\epsilon\in[0,\alpha^\star]$ with $\omega(\epsilon)\le t$ has $\epsilon<\epsilon_0$, and then, if $\epsilon>0$,
$\exp(-c'/\epsilon)\le t$, that is, $\epsilon\le c'/\log(1/t)$; so $\omega^{\leftarrow}(t)\le c'/\log(1/t)$. For
$t=t_{n,m}$, $\log(1/t_{n,m})\ge\frac12\log\underline n-\frac12\log\log\underline n-\log C_t$, which exceeds
$(c'/c)\log\underline n$ for all large $\underline n$, because $c'/c<\frac12$; then also $t_{n,m}<\min\{1,\omega(\epsilon_0)\}$,
and $\omega^{\leftarrow}(t_{n,m})\le c/\log\underline n$. Part (i), in the form of part (ii) of
Theorem~\ref{thm:lorate}, gives the claim.
\end{proof}

\subsection{Sobolev kernels}

Let $\|f\|_{H^s}^2=(2\pi)^{-d}\int|\widehat f(\xi)|^2(1+|\xi|^2)^s\,d\xi$. If
$\widehat{k_0}(\xi)\ge c(1+|\xi|^2)^{-s}$, then by \eqref{eq:newnative} every continuous $f\in L^2$ with
$\|f\|_{H^s}<\infty$ lies in $\Hh$, with $\|f\|_\Hh^2\le\|f\|_{H^s}^2/c$. Such a kernel is characteristic, as
$\widehat{k_0}>0$ \citep{sriperumbudur2010hilbert}. The approximant is a mollified indicator. For a set $E$ we use the
translation bound $\mathrm{Leb}(E\,\triangle\,(E+h))\le\Pi|h|$ for all $h\in\R^d$, which holds with $\Pi$ the perimeter of $E$
\citep[Chapter~5]{evans2015measure}, and, in $\R$, with $\Pi=2N$ for a union of $N$ disjoint intervals.

\begin{lemma}[Norm of a mollified indicator]
\label{lem:newsobnorm}
Let $s>1/2$, let $E\subset\R^d$ be bounded and measurable with $\mathrm{Leb}(E\,\triangle\,(E+h))\le\Pi|h|$ for all $h$, let
$b\in C_c^\infty$ be a probability density supported in the closed unit ball, $b_w=w^{-d}b(\cdot/w)$, and
$f_w=\mathbf 1_E*b_w$. Let $m$ be the smallest integer with $4m>2s-1$, $C_b=\|\Delta^mb\|_{L^1}$,
$j_d(y)$ the mean of $\cos(yv_1)$ for $v$ uniform in the unit ball, and $c_d=1-\sup_{y\ge1}j_d(y)$. Then $c_d>0$,
$0\le f_w\le1$, $f_w\in C_c^\infty$, and for $0<w\le1$
\[
\|f_w\|_{H^s}^2\ \le\ 2^s\,\mathrm{Leb}(E)+C_{d,s,b}\,\Pi\,w^{1-2s},
\qquad
C_{d,s,b}=\frac{2^{3s}}{2c_d}\Big[\frac1{1-2^{1-2s}}+\frac{C_b^2}{1-2^{-(4m+1-2s)}}\Big].
\]
\end{lemma}

\begin{proof}
The indicator has a Fourier tail of order $1/t$ beyond frequency $t$, because translating $E$ by $h$ changes it on a set of
measure $\Pi|h|$, and mollification at width $w$ cuts the frequencies beyond $1/w$. Write $g=\mathbf 1_E$, so that
$\widehat f_w(\xi)=\widehat g(\xi)\widehat b(w\xi)$. As $b$ is a probability density, $|\widehat b|\le1$, and
$(\Delta^mb)^\wedge(\xi)=(-|\xi|^2)^m\widehat b(\xi)$ gives $|\widehat b(\xi)|\le\min(1,C_b|\xi|^{-2m})$.

We first bound the Fourier tail of $g$. For $h\in\R^d$, Plancherel's identity gives
\[
(2\pi)^{-d}\int|\widehat g(\xi)|^2|e^{ih\cdot\xi}-1|^2\,d\xi=\|g(\cdot+h)-g\|_{L^2}^2=\mathrm{Leb}(E\,\triangle\,(E-h))\le\Pi|h| ,
\]
since $g(\cdot+h)-g$ takes values in $\{-1,0,1\}$. To isolate the high frequencies we average over $h$ uniform in the
ball of radius $1/t$. The mean of $|e^{ih\cdot\xi}-1|^2=2-2\cos(h\cdot\xi)$ is $2\{1-j_d(|\xi|/t)\}$ by rotation
invariance, and the mean of $|h|$ is at most $1/t$. The function $j_d$ is continuous, smaller than $1$ on $(0,\infty)$,
since $\cos(yv_1)<1$ off a null set, and tends to $0$ at infinity by the Riemann--Lebesgue lemma, so $c_d>0$; in $\R$,
$j_1(y)=\sin y/y$ and $c_1=1-\sin1$. As $1-j_d\ge0$ everywhere and $1-j_d(|\xi|/t)\ge c_d$ for $|\xi|\ge t$, the
averaged identity gives
\[
\Upsilon(t)\coloneqq(2\pi)^{-d}\int_{|\xi|\ge t}|\widehat g(\xi)|^2\,d\xi\ \le\ \frac{\Pi}{2c_d\,t},\qquad t>0 .
\]
We now split the frequencies into the unit ball and dyadic shells. On $\{|\xi|<1\}$, $(1+|\xi|^2)^s\le2^s$ and
$|\widehat b|\le1$, which contributes at most $2^s(2\pi)^{-d}\int|\widehat g|^2=2^s\mathrm{Leb}(E)$, the first term of
the bound. On the shell $S_j=\{2^j\le|\xi|<2^{j+1}\}$, $j\ge0$, we have $(1+|\xi|^2)^s\le(2|\xi|^2)^s\le2^{3s}4^{sj}$,
$|\widehat b(w\xi)|^2\le\min\{1,C_b^2(w2^j)^{-4m}\}$ and $(2\pi)^{-d}\int_{S_j}|\widehat g|^2\le\Upsilon(2^j)\le\Pi2^{-j}/(2c_d)$,
so the shells contribute at most $\{2^{3s}\Pi/(2c_d)\}\sum_{j\ge0}2^{j(2s-1)}\min\{1,C_b^2w^{-4m}2^{-4mj}\}$. The sum grows geometrically up to the frequency $1/w$ and decays beyond it. Let $j_0\ge0$ be the integer with $2^{j_0}\le1/w<2^{j_0+1}$. As $2s-1>0$,
$\sum_{j\le j_0}2^{j(2s-1)}\le2^{j_0(2s-1)}/(1-2^{1-2s})\le w^{1-2s}/(1-2^{1-2s})$, and as $2s-1-4m<0$,
\[
\sum_{j>j_0}C_b^2w^{-4m}2^{j(2s-1-4m)}\le\frac{C_b^2w^{-4m}2^{(j_0+1)(2s-1-4m)}}{1-2^{2s-1-4m}}\le\frac{C_b^2w^{1-2s}}{1-2^{-(4m+1-2s)}} .
\]
Adding the parts gives the bound, the two geometric series giving the two terms of $C_{d,s,b}$. Finally $0\le f_w\le1$ because $0\le g\le1$ and $b_w$ is a probability density, and
$f_w$ is smooth with compact support because $b$ is and $E$ is bounded.
\end{proof}

The laws enter through an open set $G$ that carries the sign of $p-q$, the size of its boundary near the contact, and
the order at which $p-q$ vanishes there. We use the open balls $B_\rho$ of radius $\rho$ about a fixed centre.

\begin{theorem}[Margin for Sobolev kernels]
\label{thm:newsobolev}
Let $k(x,y)=k_0(x-y)$ on $\R^d$ with $\widehat{k_0}(\xi)\ge c(1+|\xi|^2)^{-s}$, $s>d/2$, and let $P,Q$ have Lebesgue
densities $p,q$ with $0<\alpha^\star<1$. Let $G\subseteq\R^d$ be open with $p\ge q$ almost everywhere on $G$ and $p\le q$
almost everywhere off $G$, and let $\rho\ge1$, $r>0$, $\ell>0$ and $w_0\in(0,1]$ be such that
\begin{enumerate}[label=(G\arabic*),leftmargin=*]
\item $\mathrm{Leb}(E_\rho\,\triangle\,(E_\rho+h))\le\Pi_\rho|h|$ for all $h$, where $E_\rho=G\cap B_\rho$.
\item $\mathrm{Leb}(\{x\in B_\rho:\dist(x,\partial G)\le t\})\le V_\rho t$ for $0<t\le w_0$.
\item $|p(x)-q(x)|\le\ell\,\dist(x,\partial G)^r$ for almost every $x\in B_\rho$ with $\dist(x,\partial G)\le w_0$.
\end{enumerate}
Let $T(\rho)=|P-Q|(\R^d\setminus B_{\rho-1})$, and for $0<w\le w_0$ let $f_w=\mathbf 1_{E_\rho}*b_w$ as in
Lemma~\ref{lem:newsobnorm}. Then $f_w\in\Hh$, $0\le f_w\le1$, $\mathrm{err}(f_w)\le\ell V_\rho w^{r+1}+T(\rho)$, and for every
$\epsilon\in[0,\alpha^\star]$
\[
D_{\alpha^\star-\epsilon}(P,Q)\ \ge\ \bar D_{\alpha^\star-\epsilon}\ \ge\ \frac{c^{1/2}\,\{\epsilon-\ell V_\rho w^{r+1}-T(\rho)\}_+}
{\{2^s\mathrm{Leb}(B_\rho)+C_{d,s,b}\Pi_\rho w^{1-2s}\}^{1/2}} .
\]
\end{theorem}

\begin{proof}
The function $f_w$ is the approximate witness of Proposition~\ref{prop:newapproxgen}(ii), so we need its norm and its
error. Lemma~\ref{lem:newsobnorm} and \eqref{eq:newnative} give $f_w\in\Hh$, $0\le f_w\le1$ and
$\|f_w\|_\Hh^2\le c^{-1}\{2^s\mathrm{Leb}(B_\rho)+C_{d,s,b}\Pi_\rho w^{1-2s}\}$, as $E_\rho\subseteq B_\rho$. For the error,
since $\mathbf 1_G=\mathbf 1_A$ almost everywhere
for $|P-Q|$, $\mathrm{err}(f_w)=\int|\mathbf 1_G-f_w|\,d|P-Q|$. Off $B_{\rho-1}$ we use $|\mathbf 1_G-f_w|\le1$, which contributes at
most $T(\rho)$. Inside $B_{\rho-1}$ the mollified indicator differs from $\mathbf 1_G$ only near $\partial G$. Let
$x\in B_{\rho-1}$. The density $b_w$ is supported in the closed ball of radius $w\le1$ about $0$, and the
closed ball $\bar B(x,w)$ lies in $B_\rho$, so $f_w(x)=\int\mathbf 1_G(x-y)b_w(y)\,dy$. If $\dist(x,\partial G)>w$, the connected
set $\bar B(x,w)$ misses $\partial G$, so it lies in the open set $G$ or in the open set $\R^d\setminus\bar G$, and
$f_w(x)=\mathbf 1_G(x)$. Hence $|\mathbf 1_G-f_w|\le\mathbf 1\{\dist(\cdot,\partial G)\le w\}$ on $B_{\rho-1}$, and (G2) and (G3) give
$\int_{B_{\rho-1}}|\mathbf 1_G-f_w|\,|p-q|\le\ell w^r\,\mathrm{Leb}(\{x\in B_\rho:\dist(x,\partial G)\le w\})\le\ell V_\rho w^{r+1}$.
Proposition~\ref{prop:newapproxgen}(ii) finishes the proof.
\end{proof}

\begin{corollary}[Exponent of the Sobolev margin]
\label{cor:newsobolev}
In the setting of Theorem~\ref{thm:newsobolev} put $\gamma=1+(2s-1)/\{2(r+1)\}$. Suppose either that $T(\rho)=0$, that
is, $|P-Q|$ vanishes off $B_{\rho-1}$, or that $d=1$, $G=(x_0,\infty)$, the balls are centered at $x_0$, (G3) holds for
every $\rho$ with the same $\ell$ and $w_0$, and $T(\rho)\le C_Te^{-\rho/a}$ for some $C_T,a>0$ and all $\rho\ge1$. Then there
are $c_1>0$ and $\epsilon_1>0$ with $\bar D_{\alpha^\star-\epsilon}\ge c_1\epsilon^\gamma$ for $\epsilon\le\epsilon_1$, and
Assumption~\ref{ass:margin} holds with this $\gamma$. Hence, with $\widehat\alpha$ computed as in
Theorem~\ref{thm:loratemain} on a grid of $J$ levels fixed in advance with mesh $\Delta_{\mathcal G}$ below
$\alpha^\star$, $\widehat\alpha\ge\alpha^\star-\Delta_{\mathcal G}-O_\Pp((n\wedge m)^{-1/(2\gamma)})$
for fixed $J$, with $1/(2\gamma)=(r+1)/\{2(r+1)+2s-1\}$, equal to $2/(2s+3)$ for $r=1$.
\end{corollary}

\begin{proof}
The width $w$ is chosen so that the error of $f_w$ is half of $\epsilon$, which balances error and norm in
Theorem~\ref{thm:newsobolev}. If $T(\rho)=0$, write $V$ and $\Pi$ for $V_\rho$ and $\Pi_\rho$, put
$w(\epsilon)=\{\epsilon/(2\ell V)\}^{1/(r+1)}$, so that $\ell Vw^{r+1}=\epsilon/2$, and let $w_1\in(0,w_0]$ be such that
$C_{d,s,b}\Pi w^{1-2s}\ge2^s\mathrm{Leb}(B_\rho)$ for $w\le w_1$, so that the second term of the denominator dominates. For
$\epsilon\le\epsilon_1=2\ell Vw_1^{r+1}$, Theorem~\ref{thm:newsobolev} gives
$\bar D_{\alpha^\star-\epsilon}\ge c^{1/2}(\epsilon/2)(2C_{d,s,b}\Pi)^{-1/2}w(\epsilon)^{s-1/2}=c_1\epsilon^\gamma$, the exponent being $1+(s-\tfrac12)/(r+1)=\gamma$. In the second case,
$E_\rho=(x_0,x_0+\rho)$, so $\Pi_\rho=2$, and $\{x:\dist(x,\partial G)\le t\}=[x_0-t,x_0+t]$, so $V_\rho=2$. Take
$\rho(\epsilon)=\max\{2,1+a\log(4C_T/\epsilon)\}$, so that $T(\rho(\epsilon))\le\epsilon/4$, and $w(\epsilon)=\{\epsilon/(8\ell)\}^{1/(r+1)}$, so
that $2\ell w^{r+1}=\epsilon/4$. The denominator in Theorem~\ref{thm:newsobolev} is
$\{2^{s+1}\rho(\epsilon)+2C_{d,s,b}w(\epsilon)^{1-2s}\}^{1/2}$, in which the second term, a power of $1/\epsilon$, dominates the
first, of order $\log(1/\epsilon)$, for small $\epsilon$. This gives $c_1$ and $\epsilon_1$ as before. In both cases, for
$\epsilon\in[\epsilon_1,\alpha^\star]$ the path is nonincreasing, so $D_{\alpha^\star-\epsilon}\ge \bar D_{\alpha^\star-\epsilon_1}\ge
c_1(\epsilon_1/\alpha^\star)^\gamma\epsilon^\gamma$, and Assumption~\ref{ass:margin} holds with $c_0=c_1(\epsilon_1/\alpha^\star)^\gamma$, since
$D\ge \bar D$. The rate is Theorem~\ref{thm:lorate}, whose part (i) alone, with parts (ii) and (iii) for the size of
$t_{n,m}$, gives the bound with probability at least $1-\eta$.
\end{proof}

A transversal crossing of a compact $C^2$ hypersurface $\partial G$, with $p-q$ continuously differentiable near it and
$|P-Q|$ compactly supported, satisfies (G1)--(G3) with $r=1$. Indeed $\Pi_\rho$ is at most the sum of the areas of $\partial G$
and $\partial B_\rho$, by the submodularity of the perimeter, $V_\rho$ is finite by the tube formula, and $\ell$ can be
taken as a Lipschitz constant of $p-q$ near $\partial G$.

\subsection{The Gaussian kernel}

For $k(x,y)=\exp\{-(x-y)^2/(2\sigma^2)\}$ on $\R$, $\widehat{k_0}(\xi)=(2\pi)^{1/2}\sigma e^{-\sigma^2\xi^2/2}$, so by
\eqref{eq:newnative} and Plancherel's identity every continuous $f\in L^2$ whose Fourier transform vanishes outside
$[-\xi_0,\xi_0]$ lies in $\Hh$, with
\begin{equation}
\|f\|_\Hh^2\le\frac{e^{\sigma^2\xi_0^2/2}}{(2\pi)^{1/2}\sigma}\,\|f\|_{L^2}^2 ,
\label{eq:newbandlim}
\end{equation}
the elementary half of the description of this RKHS by \citet{steinwart2006explicit}. Let $b$ be an even probability
density with $\widehat b=0$ outside $[-1,1]$ and $m_2=\int t^2b(t)\,dt<\infty$, let $\bar b(t)=\int_t^\infty b$ for $t\ge0$, so that
$\bar b\le1/2$, and $b_{\xi_0}(t)=\xi_0 b(\xi_0 t)$. The density $b(t)=4\pi\cos^2(t/2)/(\pi^2-t^2)^2$ qualifies: it is $h^2/\|h\|_{L^2}^2$
for $h(t)=\cos(t/2)/(\pi^2-t^2)$, whose Fourier transform is $\cos(\pi\xi)$ on $[-1/2,1/2]$ and $0$ elsewhere, so that
$\widehat b$, a multiple of the autocorrelation of $\widehat h$, vanishes outside $[-1,1]$. Moreover $\|h\|_{L^2}^2=1/(4\pi)$ and
$m_2=\int|\widehat h'|^2/\int|\widehat h|^2=\pi^2$. For $t\ge2^{1/2}\pi$, $t^2-\pi^2\ge t^2/2$, so $b(t)\le16\pi t^{-4}$ and
$\bar b(t)\le16\pi/(3t^3)$.

The laws cross once, at $x_0$: $p\ge q$ almost everywhere on $(x_0,\infty)$ and $p\le q$ almost everywhere on
$(-\infty,x_0)$, so that $A=(x_0,\infty)$ up to a null set of $|P-Q|$, and $|p-q|(x)\le\ell|x-x_0|$ for almost every
$x$ with $|x-x_0|\le r_0$, for some $\ell>0$ and $r_0\in(0,\infty]$; we read $\bar b(\infty)$ as $0$. We write
$T_+(\rho)=|P-Q|((x_0+\rho,\infty))$.

\begin{lemma}[A band-limited step]
\label{lem:newgaussstep}
For $\xi_0>0$ and $\rho>0$, $f=\mathbf 1_{(x_0,x_0+\rho)}*b_{\xi_0}$ lies in $\Hh$, $0\le f\le1$,
$\|f\|_\Hh^2\le\rho\,e^{\sigma^2\xi_0^2/2}/\{(2\pi)^{1/2}\sigma\}$, and
\[
\mathrm{err}(f)\ \le\ \frac{\ell m_2}{2\xi_0^2}+2T_+(\rho/2)+2\alpha^\star\{\bar b(\xi_0 r_0)+\bar b(\xi_0\rho/2)\} .
\]
\end{lemma}

\begin{proof}
Let $g=\mathbf 1_{(x_0,x_0+\rho)}$. As $b_{\xi_0}$ is a probability density, $f$ is continuous, in $L^2$, with
values in $[0,1]$, and $\widehat f=\widehat g\,\widehat b(\cdot/\xi_0)$ vanishes outside $[-\xi_0,\xi_0]$. Since $|\widehat b|\le1$,
$\|f\|_{L^2}\le\|g\|_{L^2}=\rho^{1/2}$, and \eqref{eq:newbandlim} gives the norm bound.

For the error we compare $f$ with $g$, and $g$ with $\mathbf 1_A$: $|\mathbf 1_A-f|\le|\mathbf 1_A-g|+|g-f|$, and $\mathbf 1_A-g=\mathbf 1_{[x_0+\rho,\infty)}$ almost everywhere, which
contributes $T_+(\rho)\le T_+(\rho/2)$. Next, $|g-f|(x)\le\int|g(x)-g(x-y)|\,b_{\xi_0}(y)\,dy$. For $x_0<x<x_0+\rho$, $g(x-y)\neq g(x)$
exactly when $y\ge x-x_0$ or $y\le x-x_0-\rho$, a set of $b_{\xi_0}$-mass $\bar b(\xi_0(x-x_0))+\bar b(\xi_0(x_0+\rho-x))$. For
$x<x_0$ only $y<-|x-x_0|$ can change $g$, and for $x>x_0+\rho$ only $y>x-x_0-\rho$. Hence, for almost every $x$, $|g-f|(x)\le\bar b(\xi_0|x-x_0|)+\bar b(\xi_0|x-x_0-\rho|)$, the leakage of the mollifier across the two ends of the interval. Integrated against $|P-Q|$, it is at most
\begin{align*}
\ell\int_{|t|\le r_0}|t|\,\bar b(\xi_0|t|)\,dt+\bar b(\xi_0 r_0)\,|P-Q|(\R)
&\le\frac{2\ell}{\xi_0^2}\int_0^\infty v\,\bar b(v)\,dv+2\alpha^\star\bar b(\xi_0 r_0)\\
&=\frac{\ell m_2}{2\xi_0^2}+2\alpha^\star\bar b(\xi_0 r_0),
\end{align*}
because $\bar b$ is nonincreasing and $\int_0^\infty v\,\bar b(v)\,dv=\int_0^\infty b(y)\,y^2/2\,dy=m_2/4$. The second term is at most
$1/2$ on $(x_0+\rho/2,\infty)$, a set of $|P-Q|$-mass $T_+(\rho/2)$, and at most $\bar b(\xi_0\rho/2)$ elsewhere, where
$|x-x_0-\rho|\ge\rho/2$. Adding the contributions gives the bound.
\end{proof}

\begin{theorem}[Margin for the Gaussian kernel]
\label{thm:newgauss}
In this setting let $\theta\in(0,1)$, $\epsilon\in(0,\alpha^\star]$ and $\xi_\epsilon=\{\ell m_2/(2\theta\epsilon)\}^{1/2}$, and let
$\rho>0$ satisfy $2T_+(\rho/2)+2\alpha^\star\{\bar b(\xi_\epsilon r_0)+\bar b(\xi_\epsilon\rho/2)\}\le(1-\theta)\epsilon/2$. Then
\[
D_{\alpha^\star-\epsilon}(P,Q)\ \ge\ \bar D_{\alpha^\star-\epsilon}\ \ge\ \frac{(1-\theta)\epsilon}2\Big\{\frac{(2\pi)^{1/2}\sigma}{\rho}\Big\}^{1/2}
\exp\Big(-\frac{\sigma^2\ell m_2}{8\theta\epsilon}\Big) .
\]
If $T_+(t)\le C_Te^{-t/a}$, then $\rho_\epsilon=\max\{2,2a\log(8C_T/((1-\theta)\epsilon))\}$ satisfies the condition for all small
$\epsilon$ when $b$ is the density above, and $\log\{1/\bar D_{\alpha^\star-\epsilon}\}\le\sigma^2\ell m_2/(8\theta\epsilon)+\log(1/\epsilon)+O(\log\log(1/\epsilon))$.
\end{theorem}

\begin{proof}
The bandwidth $\xi_\epsilon$ is chosen so that the smoothing error near $x_0$ is the fraction $\theta$ of $\epsilon$,
and the norm of the band-limited step then grows like $\exp\{\sigma^2\xi_\epsilon^2/4\}$. By
Lemma~\ref{lem:newgaussstep} with $\xi_0=\xi_\epsilon$, $\ell m_2/(2\xi_\epsilon^2)=\theta\epsilon$, so
$\mathrm{err}(f)\le\theta\epsilon+(1-\theta)\epsilon/2$ and $\epsilon-\mathrm{err}(f)\ge(1-\theta)\epsilon/2$, while
$\|f\|_\Hh\le\{\rho/((2\pi)^{1/2}\sigma)\}^{1/2}\exp\{\sigma^2\ell m_2/(8\theta\epsilon)\}$, and Proposition~\ref{prop:newapproxgen}(ii) gives
the bound. With the density above, $\bar b(\xi_\epsilon t)\le16\pi/\{3(\xi_\epsilon t)^3\}$ once $\xi_\epsilon t\ge2^{1/2}\pi$, which is
of order $\epsilon^{3/2}$ for $t=r_0$ and for $t=\rho_\epsilon/2\ge1$, so the two terms in $\bar b$ are at most $(1-\theta)\epsilon/4$
for small $\epsilon$, and $2T_+(\rho_\epsilon/2)\le2C_T\exp\{-\log(8C_T/((1-\theta)\epsilon))\}=(1-\theta)\epsilon/4$.
\end{proof}

\begin{proof}[Proof of Proposition~\ref{prop:newgauss}]
The hypotheses of Proposition~\ref{prop:newsobolev} give the setting of Theorem~\ref{thm:newgauss}, with exponential tails,
for every $\ell>\ell_0$ and a small enough $r_0$: $(p-q)(x)/(x-x_0)$ is continuous and nonnegative near $x_0$, away from $x_0$,
and tends to $\ell_0$ at $x_0$. Given $c>\pi^2\sigma^2\ell_0/8$, choose $\theta<1$ and $\ell>\ell_0$ with $\pi^2\sigma^2\ell/(8\theta)<c$. Then the
bound of Theorem~\ref{thm:newgauss}, with $m_2=\pi^2$, exceeds $\exp(-c/\epsilon)$ for small $\epsilon$. The upper bound is
Proposition~\ref{prop:newmoments}(b) below with $2j$ vanishing moments, which gives the order $\epsilon^{1+j}$, as
$M_{4j}=(4j-1)!!\,\sigma^{-4j}<\infty$ for every $j$.
\end{proof}

\begin{proof}[Proof of Proposition~\ref{prop:newrates}(iii)]
Given $c>\pi^2\sigma^2\ell_0/4$, take $c'\in(\pi^2\sigma^2\ell_0/8,c/2)$. By the proof of Proposition~\ref{prop:newgauss},
the bound of Theorem~\ref{thm:newgauss}, with $m_2=\pi^2$, $\rho=\rho_\epsilon$ and suitable $\theta<1$ and
$\ell>\ell_0$, exceeds $\exp(-c'/\epsilon)$ for $0<\epsilon\le\epsilon_0$, for some $\epsilon_0\in(0,\alpha^\star]$. The
proof of Theorem~\ref{thm:newgauss} obtains this bound as a lower bound on $(\epsilon-\mathrm{err}(f))/\|f\|_\Hh$ for a
function $f\in\mathcal F$, so it is also a lower bound on $\omega_\Hh(\epsilon)$. By
Proposition~\ref{prop:newapproxgen}(iii), $\omega_\Hh$ is nondecreasing and at most
$\bar D_{\alpha^\star-\epsilon}=(1-\alpha^\star+\epsilon)D_{\alpha^\star-\epsilon}$, so part (ii) applies to $\omega=\omega_\Hh$,
and $c>2c'$.
\end{proof}

\subsection{Explicit trimmings and sharpness}

Upper bounds on the path come from explicit trimmings: nonnegative measures $P'\le P$ and $Q'\le Q$ of equal mass
$1-\alpha$ give $\bar D_{\alpha}\le\|\int\phi\,d(P'-Q')\|_\Hh$. If $\nu=P'-Q'$ has total variation $|\nu|$, support in $[x_0-W,x_0+W]$
and $\int(x-x_0)^j\,d\nu=0$ for $0\le j<N$, where $N\ge1$ and $j=0$ holds because the masses agree, then
$|e^{it}-\sum_{j<N}(it)^j/j!|\le|t|^N/N!$ gives
\begin{equation}
|\widehat\nu(\xi)|\le|\nu|\min\{1,(|\xi|W)^N/N!\} .
\label{eq:newmomentdecay}
\end{equation}

\begin{proposition}[Sobolev order below $3/2$]
\label{prop:newupperone}
Let $d=1$, $k(x,y)=k_0(x-y)$ with $\widehat{k_0}(\xi)\le C|\xi|^{-2s}$ for some $s\in(1/2,3/2)$, and suppose that
$p-q\ge\ell'(x-x_0)^r$ almost everywhere on $(x_0,x_0+r_0)$ and $q-p\ge\ell'(x_0-x)^r$ almost everywhere on $(x_0-r_0,x_0)$,
for some $\ell',r,r_0>0$. Then for $0<\epsilon\le\ell'r_0^{r+1}/(r+1)$, with $W_\epsilon=\{(r+1)\epsilon/\ell'\}^{1/(r+1)}$,
\[
\bar D_{\alpha^\star-\epsilon}\ \le\ 2\epsilon\,W_\epsilon^{s-1/2}\Big\{\frac C{2\pi}\Big(\frac2{3-2s}+\frac2{2s-1}\Big)\Big\}^{1/2} .
\]
\end{proposition}

\begin{proof}
We keep the common part together with the excess of $P$ on a short interval to the right of $x_0$ and that of $Q$ on
one to the left, so that the two trimmed parts differ by a measure of mass zero carried near $x_0$, whose Fourier
transform vanishes at the origin. The map $w\mapsto\int_{x_0}^{x_0+w}(p-q)$ is continuous, nondecreasing and at least $\ell'w^{r+1}/(r+1)$ on $[0,r_0]$, so some
$w_P\le W_\epsilon$ gives it the value $\epsilon$, and likewise some $w_Q\le W_\epsilon$ on the left. As $p-q=(p-q)_+$ almost everywhere
on $(x_0,x_0+w_P)$, the measure $P'=P\wedge Q+(p-q)\mathbf 1_{(x_0,x_0+w_P)}\,dx$ satisfies $P'\le P\wedge Q+(P-P\wedge Q)=P$ and has mass
$1-\alpha^\star+\epsilon$. The measure $Q'$ is built in the same way, and $\nu=P'-Q'=(p-q)\mathbf 1_{(x_0-w_Q,x_0+w_P)}\,dx$ has $|\nu|=2\epsilon$, mass
$0$ and support in $[x_0-W_\epsilon,x_0+W_\epsilon]$. By \eqref{eq:newbochner} and \eqref{eq:newmomentdecay} with $N=1$,
$\|\int\phi\,d\nu\|_\Hh^2\le(2\pi)^{-1}C(2\epsilon)^2\int\min(1,\xi^2W_\epsilon^2)|\xi|^{-2s}\,d\xi$, and the integral is
$W_\epsilon^{2s-1}\{2/(3-2s)+2/(2s-1)\}$.
\end{proof}

For larger $s$, and for the Gaussian kernel, the difference of the trimmings must have more vanishing moments. They
are created by removing from the common part a small polynomial density near $x_0$, which needs the core to have
positive density there.

\begin{proposition}[Trimmings with vanishing moments]
\label{prop:newmoments}
Let $d=1$, let $p,q$ be continuous on $(x_0-r_0,x_0+r_0)$, let $e=p-q$ be differentiable at $x_0$ with
$e'(x_0)=\ell_0>0$, and let $p(x_0)=q(x_0)>0$. For every integer $N\ge1$ there are $T_N\ge1$ and $w_N>0$ such that for
$0<w\le w_N$ there are nonnegative measures $P'\le P$ and $Q'\le Q$ of equal mass $1-\alpha^\star+\epsilon(w)$, with
$\epsilon(w)\ge\ell_0w^2/8$, whose difference $\nu$ satisfies $|\nu|\le2\ell_0w^2$, $\supp\nu\subseteq[x_0-T_Nw,x_0+T_Nw]$ and
$\int(x-x_0)^j\,d\nu=0$ for $0\le j<N$. Consequently, with $w=(8\epsilon/\ell_0)^{1/2}$ and $0<\epsilon\le\ell_0w_N^2/8$:
\begin{enumerate}[label=(\alph*),leftmargin=*]
\item if $\widehat{k_0}(\xi)\le C(1+\xi^2)^{-s}$ with $s>1/2$ and $N>s-1/2$, then
$\bar D_{\alpha^\star-\epsilon}\le2\ell_0w^2(T_Nw)^{s-1/2}\{CI_{N,s}/(2\pi)\}^{1/2}$, which is at most a constant times
$\epsilon^{(2s+3)/4}$, where $I_{N,s}=\int\min(1,v^{2N}/N!^2)|v|^{-2s}\,dv<\infty$.
\item if $M_{2N}=(2\pi)^{-1}\int\xi^{2N}\widehat{k_0}(\xi)\,d\xi<\infty$, then
$\bar D_{\alpha^\star-\epsilon}\le2\ell_0w^2(T_Nw)^NM_{2N}^{1/2}/N!$, which is at most a constant times $\epsilon^{1+N/2}$. For the
Gaussian kernel $M_{2N}=(2N-1)!!\,\sigma^{-2N}$.
\end{enumerate}
\end{proposition}

\begin{proof}
We add the excess $e$ on $(x_0-w,x_0+w)$ to the common part and cancel its first $N$ moments with a polynomial density
removed from the common part on a fixed multiple of that interval. For $|t|<1$ let $F_w(t)=e(x_0+wt)/(\ell_0w)$, and
$F_w(t)=0$ otherwise. Differentiability at $x_0$ gives $\sup_t|F_w(t)-t\mathbf 1\{|t|<1\}|\to0$ as $w\to0$.

Let $T\ge1$ and let $\mathcal G_T$ map $M\in\R^N$ to the polynomial $G$ of degree less than $N$ with
$\int_{-T}^Tt^jG(t)\,dt=M_j$ for $0\le j<N$, which exists and is unique because the Gram matrix of the monomials on
$[-T,T]$ is positive definite. The map $\mathcal G_T$ is linear. Put $M_j(w)=-\int t^jF_w(t)\,dt$, so that $M(w)\to M(0)$,
with $M_0(0)=0$ and $|M_j(0)|=|\int_{-1}^1t^{j+1}dt|\le2/3$. For $G_0=\mathcal G_T[M(0)]$, the polynomial $v\mapsto TG_0(Tv)$ has
moments $T^{-j}M_j(0)$ on $[-1,1]$, so it is $\mathcal G_1$ applied to a vector of Euclidean norm at most
$(2/3)N^{1/2}/T$, as $M_0(0)=0$. Hence $\int_{-T}^T|G_0|=\int_{-1}^1T|G_0(Tv)|\,dv\le c_N/T$ with
$c_N=(2/3)N^{1/2}\|\mathcal G_1\|$, the operator norm from $\R^N$ to $L^1[-1,1]$. Fix $T=T_N=\max(1,8c_N)$, so that
$\int|G_0|\le1/8$, and let $B_N=\sup_{[-T,T]}|G_0|+1$.

The polynomials $G_w=\mathcal G_T[M(w)]$ converge to $G_0$ uniformly on $[-T,T]$. Choose $w_N>0$ so small that for $w\le w_N$ the following hold. First, $\int|G_w|\le1/4$
and $\sup|G_w|\le B_N$. Next, $Tw\le r_0$, $\min(p,q)\ge p(x_0)/2$ on $[x_0-Tw,x_0+Tw]$ and $\ell_0wB_N\le p(x_0)/2$. Finally,
$e>0$ on $(x_0,x_0+w)$, $e<0$ on $(x_0-w,x_0)$, and both $\int_0^we(x_0+t)\,dt$ and $-\int_{-w}^0e(x_0+t)\,dt$ lie between
$3\ell_0w^2/8$ and $5\ell_0w^2/8$, which is possible because both equal $\ell_0w^2\{1/2+o(1)\}$.

Let $a_P=e\,\mathbf 1_{(x_0,x_0+w)}$, $a_Q=-e\,\mathbf 1_{(x_0-w,x_0)}$, $g(x)=\ell_0wG_w((x-x_0)/w)\mathbf 1\{|x-x_0|\le Tw\}$, and
$P'=P\wedge Q+(a_P-g_-)\,dx$, $Q'=P\wedge Q+(a_Q-g_+)\,dx$. As $|g|\le\ell_0wB_N\le\min(p,q)$ where $g\neq0$, both measures are nonnegative,
and $P'\le P\wedge Q+(p-q)_+\,dx=P$, $Q'\le Q$. Their difference is $\nu=(e\,\mathbf 1_{(x_0-w,x_0+w)}+g)\,dx$, and for $j<N$
\begin{align*}
\int(x-x_0)^j\,d\nu&=\ell_0w^{j+2}\Big\{\int t^jF_w(t)\,dt+\int_{-T}^Tt^jG_w(t)\,dt\Big\}\\
&=\ell_0w^{j+2}\{-M_j(w)+M_j(w)\}=0 .
\end{align*}
The case $j=0$ says that the masses agree, and they equal $1-\alpha^\star+\epsilon(w)$ with $\epsilon(w)=\int a_P-\int g_-$. As
$\int|g|=\ell_0w^2\int|G_w|\le\ell_0w^2/4$, $\epsilon(w)\ge3\ell_0w^2/8-\ell_0w^2/4=\ell_0w^2/8$ and
$|\nu|\le5\ell_0w^2/4+\ell_0w^2/4\le2\ell_0w^2$. The support of $\nu$ lies in $[x_0-Tw,x_0+Tw]$ because $T\ge1$.

For the bounds, $w=(8\epsilon/\ell_0)^{1/2}\le w_N$ gives $\epsilon\le\epsilon(w)$, and $\bar D$ is nonincreasing, so
$\bar D_{\alpha^\star-\epsilon}\le \bar D_{\alpha^\star-\epsilon(w)}\le\|\int\phi\,d\nu\|_\Hh$. By \eqref{eq:newbochner} and \eqref{eq:newmomentdecay}
with $W=Tw$, in case (a) $\|\int\phi\,d\nu\|_\Hh^2\le(2\pi)^{-1}C|\nu|^2\int\min\{1,(|\xi|Tw)^{2N}/N!^2\}|\xi|^{-2s}\,d\xi
=(2\pi)^{-1}C|\nu|^2(Tw)^{2s-1}I_{N,s}$, where $I_{N,s}$ is finite because $2N-2s>-1$ near $0$ and $2s>1$ at infinity. In case
(b), $\|\int\phi\,d\nu\|_\Hh^2\le|\nu|^2(Tw)^{2N}M_{2N}/N!^2$. With $|\nu|\le2\ell_0w^2$ and $w^2=8\epsilon/\ell_0$ these are the
bounds. For the Gaussian kernel, $M_{2N}=(-1)^Nk_0^{(2N)}(0)=(2N-1)!!\,\sigma^{-2N}$.
\end{proof}

\begin{proof}[Proof of Proposition~\ref{prop:newsobolev}]
For the lower bound, $(p-q)(x)/(x-x_0)$ is continuous near $x_0$ with limit $\ell_0>0$, so for some $w_0\le1$ and $\ell<\infty$,
$|p-q|(x)\le\ell|x-x_0|$ for $|x-x_0|\le w_0$, and Corollary~\ref{cor:newsobolev}, second case, applies with $r=1$, since
$\widehat{k_0}\ge c(1+\xi^2)^{-s}$ with $s>1/2=d/2$. For the upper bound, Proposition~\ref{prop:newmoments}(a) with any
$N>s-1/2$ applies, and $D_{\alpha^\star-\epsilon}\le \bar D_{\alpha^\star-\epsilon}/(1-\alpha^\star)$. For $s<3/2$ and $\alpha^\star<1$,
Proposition~\ref{prop:newupperone} with $r=1$ already gives the upper bound, without $p(x_0)>0$.
\end{proof}

\subsection{The converse}

For a version of the analytic baseline run with the larger threshold $\tau^{\mathrm S}$ of \eqref{eq:tauA} alone, an
upper bound on the path gives a lower bound on the error. For $\tau^{\mathrm P}$, the threshold the baseline uses, no
such converse is known.

\begin{proposition}[The margin governs the rate]
\label{prop:newconverse}
Let (A1) and (A2) of Assumption~\ref{ass:regularity} hold, let $\alpha^\star<1$ and $0<\delta<\delta'<1$, let
$\widehat\alpha_{\mathrm{an}}^{\mathrm S}$ be \eqref{eq:alphaest} computed with the threshold $\tau^{\mathrm S}(\cdot,\delta)$, and put
$g_{n,m}=\Delta_\phi\{(1/n+1/m)/2\}^{1/2}\{\log^{1/2}(1/\delta)-\log^{1/2}(1/\delta')\}$. Then, with probability at
least $1-\delta'$,
\[
\alpha^\star-\widehat\alpha_{\mathrm{an}}^{\mathrm S}\ \ge\ \sup\big\{\epsilon\in[0,\alpha^\star]:\ D_{\alpha^\star-\epsilon}(P,Q)\le g_{n,m}\big\} ,
\]
a set that contains $0$, since the path vanishes at $\alpha^\star$.
\end{proposition}

\begin{proof}[Proof of Proposition~\ref{prop:newconverse}]
On the event where the statistic exceeds the path by at most $\tau^{\mathrm S}(\cdot,\delta')$, every level at which
the path lies below the gap between the thresholds at $\delta$ and at $\delta'$ is accepted. Let $E^+$ be the event of Theorem~\ref{thm:twosided} at level $\delta'$, of probability at least $1-\delta'$. On $E^+$, for every
$\epsilon\in[0,\alpha^\star]$, $\widehat D_{\alpha^\star-\epsilon}\le D_{\alpha^\star-\epsilon}+\tau^{\mathrm S}_{n,m}(\alpha^\star-\epsilon,\delta')$. With
$\beta=(1-\alpha^\star+\epsilon)^{-1}\ge1$,
\begin{align*}
&\tau^{\mathrm S}_{n,m}(\alpha^\star-\epsilon,\delta)-\tau^{\mathrm S}_{n,m}(\alpha^\star-\epsilon,\delta')\\
&\qquad=\beta\Delta_\phi\{(1/n+1/m)/2\}^{1/2}\{\log^{1/2}(1/\delta)-\log^{1/2}(1/\delta')\}\ge g_{n,m} .
\end{align*}
So if $D_{\alpha^\star-\epsilon}\le g_{n,m}$, then $\widehat D_{\alpha^\star-\epsilon}\le\tau^{\mathrm S}_{n,m}(\alpha^\star-\epsilon,\delta)$, the level
$\alpha^\star-\epsilon$ is accepted, and $\widehat\alpha_{\mathrm{an}}^{\mathrm S}\le\alpha^\star-\epsilon$ by the definition \eqref{eq:alphaest}. The event does
not depend on $\epsilon$, so this holds for all such $\epsilon$ at once.
\end{proof}

Under the hypotheses of Proposition~\ref{prop:newsobolev}, let $\epsilon_0>0$ be such that
$D_{\alpha^\star-\epsilon}\le c_2\epsilon^{(2s+3)/4}$ for $\epsilon\le\epsilon_0$; then on $E^+$,
$\alpha^\star-\widehat\alpha_{\mathrm{an}}^{\mathrm S}\ge\min\{\alpha^\star,\epsilon_0,(g_{n,m}/c_2)^{4/(2s+3)}\}$, of order $(n\wedge m)^{-2/(2s+3)}$ since
$1/n+1/m\ge1/(n\wedge m)$. Under the hypotheses of Proposition~\ref{prop:newgauss}, the bound
$D_{\alpha^\star-\epsilon}\le C_j\epsilon^j$ gives $\alpha^\star-\widehat\alpha_{\mathrm{an}}^{\mathrm S}\ge c(n\wedge m)^{-1/(2j)}$ on $E^+$ for every
$j$ and $n\wedge m$ large, so the error is not $O((n\wedge m)^{-a})$ for any $a>0$.

\subsection{The excess of the population ramp}
\label{sec:excessproofs}

This subsection proves Proposition~\ref{prop:excess}. Its arguments use $\alpha^\star<1$ only in part (ii) of the
proposition: Lemmas~\ref{lem:excessbathtub} and~\ref{lem:excessid} and parts (i) and (iii) hold for any pair of
probability measures with $\alpha<\TV(P,Q)\le1$, which Section~\ref{sec:laplaceproof} uses for empirical laws. Fix
$\alpha<\alpha^\star$, with $\bar D$ and $\bar{\mathcal C}_\alpha$ as in
Section~\ref{sec:generalmargin}; $\bar D_\alpha>0$ by \eqref{eq:zeroeq}. By Lemma~\ref{lem:attainment} there are parts
$P'$ of $P$ and $Q'$ of $Q$ of mass $1-\alpha$ such that
$d=\mu_{P'}-\mu_{Q'}$ is the element of minimal norm of $\bar{\mathcal C}_\alpha(P)-\bar{\mathcal C}_\alpha(Q)$, with $\|d\|_\Hh=\bar D_\alpha$, and
$f^\star=d/\bar D_\alpha$ is the witness of Lemma~\ref{lem:popramp}, with the same cut points $t^\star$ and $u^\star$.

\begin{lemma}[Bathtub for the closest parts]
\label{lem:excessbathtub}
Every measure $P''\le P$ satisfies $\int f^\star\,dP''\ge\int f^\star\,dP'+t^\star\{P''(\mathcal X)-(1-\alpha)\}$, and
every $Q''\le Q$ satisfies $\int f^\star\,dQ''\le\int f^\star\,dQ'+u^\star\{Q''(\mathcal X)-(1-\alpha)\}$. Moreover
$P'=P$ on $\{f^\star<t^\star\}$, $P'(\{f^\star>t^\star\})=0$, $Q'=Q$ on $\{f^\star>u^\star\}$ and
$Q'(\{f^\star<u^\star\})=0$.
\end{lemma}

\begin{proof}
As $d$ has minimal norm in a convex set, $\langle d,c-d\rangle_\Hh\ge0$ for every $c$ in it. With
$c=\mu_{P''}-\mu_{Q'}$ for a part $P''$ of $P$ of mass $1-\alpha$ and the reproducing property, $P'$ minimizes
$\int f^\star\,dP''$ over these parts; likewise $Q'$ maximizes $\int f^\star\,dQ''$ over the parts of $Q$ of mass
$1-\alpha$. Suppose $(P-P')(\{f^\star<t^\star\})>0$. As $t^\star$ is a $(1-\alpha)$-quantile,
$P(\{f^\star<t^\star\})\le1-\alpha$, so $P'(\{f^\star\ge t^\star\})=(1-\alpha)-P'(\{f^\star<t^\star\})\ge
(P-P')(\{f^\star<t^\star\})>0$. Moving a small mass from $P'$ on $\{f^\star\ge t^\star\}$ to $P-P'$ on
$\{f^\star<t^\star\}$ gives another part of mass $1-\alpha$ with a smaller integral of $f^\star$, a contradiction; so
$P'=P$ on $\{f^\star<t^\star\}$. If $P'(\{f^\star>t^\star\})>0$, then, as $P(\{f^\star\le t^\star\})\ge1-\alpha$, the
measure $P-P'$ has at least that mass on $\{f^\star=t^\star\}$, and moving mass there from $\{f^\star>t^\star\}$ lowers
the integral; so $P'(\{f^\star>t^\star\})=0$. The claims on $Q'$ follow in the same way from
$Q(\{f^\star>u^\star\})\le1-\alpha\le Q(\{f^\star\ge u^\star\})$.

For the inequality, let $P''\le P$ have any mass and $\lambda=P''-P'$. On $\{f^\star<t^\star\}$, $\lambda=P''-P\le0$
and $f^\star<t^\star$; on $\{f^\star>t^\star\}$, $\lambda=P''\ge0$ and $f^\star>t^\star$. So
$\int_Af^\star\,d\lambda\ge t^\star\lambda(A)$ for each of the sets $A=\{f^\star<t^\star\}$,
$\{f^\star=t^\star\}$ and $\{f^\star>t^\star\}$, and adding gives
$\int f^\star\,d\lambda\ge t^\star\lambda(\mathcal X)=t^\star\{P''(\mathcal X)-(1-\alpha)\}$. The bound for
$Q''$ is the same argument with the inequalities reversed.
\end{proof}

\begin{lemma}[Excess identity]
\label{lem:excessid}
\[
\E_P\rp^\star-\E_Q\rp^\star-\alpha
=\frac{\bar D_\alpha+\int(u^\star-f^\star)_+\,dP+\int(f^\star-t^\star)_+\,dQ}{t^\star-u^\star} .
\]
\end{lemma}

\begin{proof}
By Lemma~\ref{lem:excessbathtub}, $P-P'$ has mass $\alpha$ and is carried by $\{f^\star\ge t^\star\}$, where
$\rp^\star=1$, and $Q-Q'$ is carried by $\{f^\star\le u^\star\}$, where $\rp^\star=0$. Hence the left side
is $\int\rp^\star\,d(P'-Q')$. With the clipped witness $\tilde f=(f^\star\vee u^\star)\wedge t^\star$,
$\rp^\star=(\tilde f-u^\star)/(t^\star-u^\star)$, and as $P'$ and $Q'$ have the same mass, the left side is
$\int\tilde f\,d(P'-Q')/(t^\star-u^\star)$. On $\{f^\star\le t^\star\}$, which carries $P'$,
$\tilde f=f^\star+(u^\star-f^\star)_+$, and on $\{f^\star\ge u^\star\}$, which carries $Q'$,
$\tilde f=f^\star-(f^\star-t^\star)_+$. So $\int\tilde f\,d(P'-Q')=\langle f^\star,d\rangle_\Hh+\int(u^\star-f^\star)_+\,dP'
+\int(f^\star-t^\star)_+\,dQ'$, the first term is $\bar D_\alpha$, and $P'$ and $Q'$ may be replaced by $P$ and $Q$ in the
other two, because $P'=P$ on $\{f^\star<u^\star\}\subseteq\{f^\star<t^\star\}$ and $Q'=Q$ on
$\{f^\star>t^\star\}\subseteq\{f^\star>u^\star\}$.
\end{proof}

The two integrals are overshoots of the witness beyond its cut points, each against the law that is not trimmed there.

\begin{proof}[Proof of Proposition~\ref{prop:excess}]
The first inequality of (i) is Lemma~\ref{lem:excessid}, the integrals being nonnegative. For the second, let
$\alpha'\in[0,1)$. As $f^\star$ has unit norm, the Cauchy--Schwarz inequality gives
$\bar D_{\alpha'}\ge\int f^\star\,dP''-\int f^\star\,dQ''$ for the closest parts $P''$ of $P$ and $Q''$ of $Q$ of mass
$1-\alpha'$, and by Lemma~\ref{lem:excessbathtub} the right side is at least
$\int f^\star\,d(P'-Q')-(t^\star-u^\star)(\alpha'-\alpha)=\bar D_\alpha-(t^\star-u^\star)(\alpha'-\alpha)$.

For (ii), let $\alpha=\alpha^\star-\epsilon$ with $\epsilon\le\epsilon_0/2$ and take $\alpha'=\alpha^\star-2\epsilon$ in
(i): $t^\star-u^\star\le\{\bar D_{\alpha^\star-2\epsilon}-\bar D_{\alpha^\star-\epsilon}\}/\epsilon\le
D_{\alpha^\star-2\epsilon}/\epsilon\le c_2(2\epsilon)^\gamma/\epsilon$, while
$\bar D_{\alpha^\star-\epsilon}\ge(1-\alpha^\star)D_{\alpha^\star-\epsilon}\ge(1-\alpha^\star)c_1\epsilon^\gamma$. The first
inequality of (i) gives the bound. Proposition~\ref{prop:newsobolev} gives the two-sided bound with
$\gamma=(2s+3)/4$.

For (iii), write $d(x)=\int k(x,y)\,d\nu(y)$ with $\nu=P'-Q'$, a finite signed measure with $\nu(\R)=0$. Since
$(1-\sigma^2\partial_x^2)e^{-|x|/\sigma}=2\sigma\delta_0$ in the sense of distributions, $d-\sigma^2d''=2\sigma\nu$. The function
$d$ is continuous and tends to $0$ at $\pm\infty$, and it is not zero, as $\|d\|_\Hh=\bar D_\alpha>0$; by Fubini's theorem
$\int d(x)\,dx=2\sigma\nu(\R)=0$, so $d$ takes positive values and attains its maximum $M>0$ at some $z_0$. Suppose that
$M>\bar D_\alpha t^\star$. The open set $O=\{f^\star>t^\star\}$ contains $z_0$, and on it $P'=0$ and $Q'=Q$ by
Lemma~\ref{lem:excessbathtub}, as $t^\star>u^\star$, so $\nu=-Q\le0$ on $O$. On an open interval $I\ni z_0$ inside $O$
on which $d>M/2$, the distribution $\sigma^2d''=d-2\sigma\nu$ is at least $M/2$, so $\tilde d(x)=d(x)-M(x-z_0)^2/(4\sigma^2)$ is
convex on $I$. But $\frac12\{\tilde d(z_0+h)+\tilde d(z_0-h)\}\le d(z_0)-Mh^2/(4\sigma^2)<\tilde d(z_0)$ for small $h>0$, a contradiction.
Hence $f^\star\le t^\star$ on $\R$, and the same argument at the minimum of $d$, where $\nu=P\ge0$ on
$\{f^\star<u^\star\}$, gives $f^\star\ge u^\star$. Both overshoots in Lemma~\ref{lem:excessid} vanish.
\end{proof}

In (iii) the witness is therefore flat on the set it trims from each law, $P(\{f^\star=t^\star\})\ge\alpha$ and
$Q(\{f^\star=u^\star\})\ge\alpha$. The argument rests on $k$ being the Green function of a second-order operator with a
maximum principle.

\subsection{Proof of Theorem~\ref{thm:laplacerate}}
\label{sec:laplaceproof}

Theorem~\ref{thm:laplacerate} is the following theorem with its constants absorbed in the $O$: as $\eta<1$,
$L_{n,m}\ge\log(n\vee m)$, so $(\log(n\vee m))^{1/4}L_{n,m}^{1/2}\le L_{n,m}^{3/4}$ and
$(\log(n\vee m))^{1/2}L_{n,m}\le L_{n,m}^{3/2}$.

\begin{theorem}[Theorem~\ref{thm:laplacerate} with its constants]
\label{thm:laplaceconst}
Let $P,Q$ satisfy the hypotheses of Proposition~\ref{prop:newsobolev}, let $k(x,y)=e^{-|x-y|/\sigma}$ with $\sigma>0$
fixed in advance, and let $n\wedge m\ge10$. Compute $\widehat\alpha$ at level $\delta$ as in
Proposition~\ref{prop:loconsistency}, with the batches and fractions of Section~\ref{sec:betting},
$\lambda_{\max}=\tfrac12$, initial variance $v_0^2\le4$ in \eqref{eq:fractions} and the running variance floored at
$v^2_{\min}\in(0,4]$, an exact minimizer of \eqref{eq:qp} at each level, and a grid of $J$ levels in $(0,1)$ fixed
in advance whose mesh below $\alpha^\star$ is $\Delta_{\mathcal G}$, but with cut points that are genuine empirical
quantiles: any $(1-\alpha)$-quantile of the witness values at the $X$'s and any $\alpha$-quantile of those at the
$Y$'s of the fitting batches, computed with or without the own terms. Let $\eta\in(0,1)$, $L_\delta=\log(1/\delta)$
and $L_{n,m}=\log\{(n\vee m)/\eta\}$.
\begin{enumerate}[label=(\roman*),leftmargin=*]
\item If \eqref{eq:condR} holds for some $r_\ast\in(0,1]$ and $c_\ast>0$, then with probability at least $1-\delta-\eta$,
\[
\alpha^\star-\Delta_{\mathcal G}-C\,\frac{(\log(n\vee m))^{1/4}L_{n,m}^{1/2}
+v_{\min}^{-2}\{L_\delta^{1/2}+(\log(2J/\eta))^{1/2}\}}{(n\wedge m)^{1/2}}
\ \le\ \widehat\alpha\ \le\ \alpha^\star .
\]
\item If instead $(P+Q)(\{x:|x-x_0|>t\})\le C'e^{-t/C'}$ for all $t>0$, then with probability at least
$1-\delta-\eta$,
\[
\alpha^\star-\Delta_{\mathcal G}-C\bigg[\bigg\{\frac{(\log(n\vee m))^{1/2}L_{n,m}}{n\wedge m}\bigg\}^{4/9}
+\frac{L_\delta^{1/2}+(\log(2J/\eta))^{1/2}}{v_{\min}^{2}(n\wedge m)^{1/2}}\bigg]
\le\widehat\alpha\le\alpha^\star .
\]
\end{enumerate}
Here $C$ depends only on $\sigma$ and on the pair through $\alpha^\star$, the constants of
Proposition~\ref{prop:newsobolev} with $s=1$, the tail constant $C'$ (that of $P+Q$ in (ii)), and, in (i), $c_\ast$, $r_\ast$ and
the supremum of $p+q$ on $[x_0-r_\ast,x_0+r_\ast]$, which is finite once $r_\ast$ is small enough.
\end{theorem}

The constant $C$ degenerates when $c_\ast$ is small, as for narrow contaminations on a wide core, and it involves the
constants $c_1$, $c_2$, $\epsilon_0$ of Proposition~\ref{prop:newsobolev}, which come from limiting arguments
(Corollary~\ref{cor:newsobolev} and Proposition~\ref{prop:newmoments}) and are not explicit; whether $C$ is uniform over
a natural class is open, and we return to it after Lemma~\ref{lem:laplocal}. The theorem needs exact minimizers of
\eqref{eq:qp}, as self-clipping (Lemma~\ref{lem:lapclip}) is a first-order property whose survival under an inexact
solve is open. In $\R^d$ the margin that Theorem~\ref{thm:newsobolev} gives a transversal crossing guarantees at best
$(n\wedge m)^{-2/(d+3)}$, as a Sobolev kernel needs $s>d/2$, and the minimax rate is slower than parametric for
$L_1$-type functionals \citep{lepski1999estimation,cai2011testing,han2020estimation}. The simulation of
Section~\ref{sec:excessmc} agrees with the rate, but its cut points are interpolated between order statistics, which
the theorem does not cover, and its grid approaches the known $\alpha^\star$ geometrically.

The proof rests on a property special to the Laplace kernel on the line: the fitted witness never leaves the band
between its cut points, so the fitted ramp maximizes the ratio of empirical excess to norm over the functions of oscillation at most
one (Lemmas~\ref{lem:lapratio} and~\ref{lem:lapclip}). Comparing this ratio with the population path at the levels
$\alpha^\star-\epsilon$ and $\alpha^\star-2\epsilon$ shows that the fitted ramp keeps a fixed fraction of $\epsilon$ as
excess, up to two fluctuation terms (Lemma~\ref{lem:lapcompare}), which we bound through the norm and the error of the
function, localized near $x_0$ by \eqref{eq:condR} (Lemmas~\ref{lem:lapintervals} to~\ref{lem:lapdev}); a fixed-point
argument closes the bound (Proposition~\ref{prop:lapfitted}), and the bets are those of Theorem~\ref{thm:lorate}. Rescaling the data by $1/\sigma$ leaves the witness values, cut points and ramps unchanged and
preserves the hypotheses with rescaled constants, so we take $\sigma=1$. The reproducing kernel Hilbert space of
$e^{-|x-y|}$ is the Sobolev space $H^1(\R)$ with $\|f\|_\Hh^2=\tfrac12\int(f^2+f'^2)$, since $e^{-|x|}$ has Fourier
transform $2/(1+\omega^2)$ \citep[Theorem~10.12]{wendland2005scattered}. Every $f\in\Hh$ is absolutely continuous,
tends to $0$ at $\pm\infty$, and satisfies
\begin{equation}
|f(x)|\le\|f\|_\Hh,\qquad\int_a^b|f'|\le\sqrt{2(b-a)}\,\|f\|_\Hh .
\label{eq:laphone}
\end{equation}
Write $\operatorname{osc}(f)=\sup f-\inf f$; for $f\in\Hh$ with $\operatorname{osc}(f)\le1$, $|f|\le1$, as $f$ tends to
$0$ at infinity. For bounded $f$ let
$L(f)=\int f\,d(P-Q)$ and $\mathrm{err}(f)=\alpha^\star-L(f)$, the error of Section~\ref{sec:linearmargin} when $f$ takes
values in $[0,1]$. If $\operatorname{osc}(f)\le1$, then $0\le\mathrm{err}(f)\le2$, and with $M_f=\sup f$, $m_f=\inf f$,
\begin{equation}
\mathrm{err}(f)=\alpha^\star\{1-(M_f-m_f)\}+\int(M_f-f)(p-q)_+\,dx+\int(f-m_f)(q-p)_+\,dx ,
\label{eq:laperr}
\end{equation}
all three terms being nonnegative: the error of a function of oscillation at most one is its loss of range plus its
distances from its maximum where $p>q$ and from its minimum where $q>p$. With $\bar D$ as in
Section~\ref{sec:generalmargin}, $\bar D_\alpha$ is the smallest $\|\mu_{P'}-\mu_{Q'}\|_\Hh$ over parts $P'\le P$,
$Q'\le Q$ of mass $1-\alpha$, and we put
$V(\epsilon)=\bar D_{\alpha^\star-\epsilon}$. For other laws $\tilde P,\tilde Q$ we write $\bar D_\alpha(\tilde P,\tilde Q)=(1-\alpha)D_\alpha(\tilde P,\tilde Q)$ for
their rescaled path and $\tilde L(f)=\int f\,d(\tilde P-\tilde Q)$.

By Proposition~\ref{prop:newsobolev} with $s=1$ there are $0<c_1\le c_2$ and $\epsilon_0\in(0,\alpha^\star]$ with
$c_1\epsilon^{5/4}\le D_{\alpha^\star-\epsilon}\le c_2\epsilon^{5/4}$ for $0<\epsilon\le\epsilon_0$. Put
\[
\rho_0=\frac{(1-\alpha^\star)c_1}{2^{5/4}c_2},\quad c_V=(1-\alpha^\star)c_1\Big(\frac{\epsilon_0}{\alpha^\star}\Big)^{5/4},
\quad\kappa_0=\min\Big\{\rho_0,\frac{(1-\alpha^\star)c_1(\epsilon_0/2)^{5/4}}{2\alpha^\star}\Big\},
\]
so that $\rho_0,\kappa_0\in(0,1)$ and $c_V>0$. Here $\rho_0$ is a lower bound on $V(\epsilon)/V(2\epsilon)$
for $\epsilon<\epsilon_0/2$, by the two-sided bound and $(1-\alpha^\star)D_{\alpha^\star-\epsilon}\le V(\epsilon)\le
D_{\alpha^\star-\epsilon}$; $c_V$ extends the lower bound to all levels, $V(\epsilon)\ge c_V\epsilon^{5/4}$ for
$\epsilon\in(0,\alpha^\star]$, as $V(\epsilon)\ge(1-\alpha^\star)D_{\alpha^\star-\epsilon}$ and $V$ is nondecreasing; and
$\kappa_0$ is the fraction of $\epsilon$ that the comparison below guarantees as excess. In (i) put $\bar p=\sup_{[x_0-r_\ast,x_0+r_\ast]}(p+q)$, which is finite once $r_\ast$ is small, as $p,q$ are continuous near
$x_0$, and shrinking $r_\ast$ keeps \eqref{eq:condR} with the same $c_\ast$; off this
interval $p+q\le|p-q|/(c_\ast r_\ast)$, so $P+Q$ inherits the exponential tails of $|P-Q|$.

\begin{lemma}[Ratio inequality]
\label{lem:lapratio}
For any probability measures $\tilde P,\tilde Q$, any $\alpha\in[0,1)$ and any $\varphi\in\Hh\setminus\{0\}$ with
$\operatorname{osc}(\varphi)\le1$, $\bar D_\alpha(\tilde P,\tilde Q)\ge\{\tilde L(\varphi)-\alpha\}/\|\varphi\|_\Hh$. In particular
$V(\epsilon')\ge\{\epsilon'-\mathrm{err}(\varphi)\}/\|\varphi\|_\Hh$ for $\epsilon'\in[0,\alpha^\star]$.
\end{lemma}

\begin{proof}
For parts $P''\le\tilde P$, $Q''\le\tilde Q$ of mass $1-\alpha$, $\int\varphi\,dP''\ge\int\varphi\,d\tilde P-\alpha\sup\varphi$ and
$\int\varphi\,dQ''\le\int\varphi\,d\tilde Q-\alpha\inf\varphi$, so $\|\varphi\|_\Hh\|\mu_{P''}-\mu_{Q''}\|_\Hh\ge
\int\varphi\,d(P''-Q'')\ge\tilde L(\varphi)-\alpha$; take the infimum over the parts.
\end{proof}

At a level $\alpha$ and a batch $b$, let $\tilde P,\tilde Q$ be the empirical laws of the fitting batches, of sizes
$n',m'$, with the pooled points pairwise distinct (almost surely, as $P$ and $Q$ have densities), let $\widehat f$ be the witness direction, $\bar{\widehat D}=\bar
D_\alpha(\tilde P,\tilde Q)$, $\widehat t=\widehat f(X)_{(\lceil n'(1-\alpha)\rceil)}$ and $\widehat u=\widehat f(Y)_{(\lfloor m'\alpha\rfloor+1)}$,
the smallest $(1-\alpha)$-quantile of the values at the $X$'s and the largest $\alpha$-quantile of those at the $Y$'s.
By Lemma~\ref{lem:lapclip}, $\widehat t$ is also the $(n'_{1-\alpha}+1)$-th smallest value at the $X$'s,
$n'_{1-\alpha}=\lfloor n'(1-\alpha)\rfloor$, the cut of Proposition~\ref{prop:bangbang}, and likewise for $\widehat u$.

\begin{lemma}[The fitted witness clips itself]
\label{lem:lapclip}
Let $\alpha\in(0,1)$. Then $\bar{\widehat D}>0$, $\widehat u\le\widehat f\le\widehat t$ on $\R$, $\widehat s=\widehat t-\widehat u
=\operatorname{osc}(\widehat f)>0$, the values of $\widehat f$ at the $X$'s of rank at least $\lceil n'(1-\alpha)\rceil$
equal $\widehat t$, and those at the $Y$'s of rank at most $\lfloor m'\alpha\rfloor+1$ equal $\widehat u$. Hence every
genuine $(1-\alpha)$-quantile of the values at the $X$'s is $\widehat t$, every genuine $\alpha$-quantile of those at the
$Y$'s is $\widehat u$, and the plain ramp is $\widehat\rp=(\widehat f-\widehat u)/\widehat s$. Moreover
$\widehat\varphi=\widehat f/\widehat s$ has $\operatorname{osc}(\widehat\varphi)=1$, $\widehat N=\|\widehat\varphi\|_\Hh=
1/\widehat s\ge\tfrac12$ and
\[
\tilde L(\widehat\varphi)-\alpha=\bar{\widehat D}\,\widehat N ,
\]
so by Lemma~\ref{lem:lapratio} $\widehat\varphi$ maximizes $\{\tilde L(\varphi)-\alpha\}/\|\varphi\|_\Hh$ over
$\operatorname{osc}(\varphi)\le1$. With the own terms removed, the values at the $X$'s of rank at least
$\lceil n'(1-\alpha)\rceil$ lie in $[\widehat t-1/(n'\bar{\widehat D}),\widehat t]$ and those at the $Y$'s of rank at most
$\lfloor m'\alpha\rfloor+1$ in $[\widehat u,\widehat u+1/(m'\bar{\widehat D})]$.
\end{lemma}

\begin{proof}
The empirical laws are themselves a pair to which the results on the population ramp apply. Since
$\TV(\tilde P,\tilde Q)=1>\alpha$, $\bar{\widehat D}>0$ by \eqref{eq:zeroeq}. As noted at the start of
Section~\ref{sec:excessproofs}, Lemmas~\ref{lem:excessbathtub} and~\ref{lem:excessid} and
Proposition~\ref{prop:excess}(iii) use only Assumption~\ref{ass:regularity}, Lemmas~\ref{lem:attainment}
and~\ref{lem:popramp}, and $\alpha<\TV(P,Q)$. They therefore hold for $\tilde P,\tilde Q$ and any
$\alpha<1=\TV(\tilde P,\tilde Q)$,
and the distributional argument of (iii) applies to the atomic measure $\nu=P''-Q''$ unchanged, as $\nu\le0$ on
$\{\widehat f>\widehat t\}$. Applied to $\tilde P,\tilde Q$ they give the two facts from which everything else is read
off: $\widehat u\le\widehat f\le\widehat t$ on $\R$, and $\tilde P\widehat\rp-\tilde Q\widehat\rp-\alpha=\bar{\widehat D}/\widehat s$. Both bounds are values of $\widehat f$ at sample points, so
$\operatorname{osc}(\widehat f)=\widehat s$, and $\widehat s>0$, as $\int\widehat f\,dx=0$ forces
$\max\widehat f>0>\min\widehat f$. The order statistics of rank at least $\lceil n'(1-\alpha)\rceil$ are at least
$\widehat t$ and at most $\max\widehat f\le\widehat t$, so they equal $\widehat t$, and likewise those of the $Y$'s of
rank at most $\lfloor m'\alpha\rfloor+1$ equal $\widehat u$. A genuine $(1-\alpha)$-quantile lies between the order
statistics of ranks $\lceil n'(1-\alpha)\rceil$ and $\lfloor n'(1-\alpha)\rfloor+1$, and a genuine $\alpha$-quantile
between those of ranks $\lceil m'\alpha\rceil\ge1$ and $\lfloor m'\alpha\rfloor+1$, so they are $\widehat t$ and
$\widehat u$. As $\widehat\rp$ and $\widehat\varphi$ differ by a constant, $\tilde L(\widehat\varphi)=\tilde
P\widehat\rp-\tilde Q\widehat\rp$, which gives the identity, and $\widehat N\ge\sup|\widehat\varphi|\ge\tfrac12$ by
\eqref{eq:laphone}, as $\operatorname{osc}(\widehat\varphi)=1$. Finally, in the normalization of \eqref{eq:witness} the
own term at $X_i$ is $w_i/\widehat D_\alpha=w_i(1-\alpha)/\bar{\widehat D}\in[0,1/(n'\bar{\widehat D})]$, and at $Y_j$ it
is at most $1/(m'\bar{\widehat D})$; lowering each value by at most $c$ lowers each order statistic by at most $c$.
\end{proof}

At the population level, let $f^\star$, $t^\star$, $u^\star$ and $\rp^\star$ be as in
Proposition~\ref{prop:excess} at $\alpha=\alpha^\star-\epsilon\in(0,\alpha^\star)$, $s^\star=t^\star-u^\star$ and
$\varphi^\star=f^\star/s^\star$. By Proposition~\ref{prop:excess}(iii), $\operatorname{osc}(\varphi^\star)\le1$,
$\|\varphi^\star\|_\Hh=1/s^\star$ and
\begin{equation}
L(\varphi^\star)-\alpha=\E_P\rp^\star-\E_Q\rp^\star-\alpha=V(\epsilon)/s^\star\le\epsilon ,
\label{eq:lappop}
\end{equation}
the last inequality because $L(\varphi^\star)\le\alpha^\star$. So $s^\star\ge V(\epsilon)/\epsilon$,
$\|\varphi^\star\|_\Hh\le c_V^{-1}\epsilon^{-1/4}$, $\mathrm{err}(\varphi^\star)\le\epsilon$, and $s^\star\le\operatorname{osc}(f^\star)\le2$.

\begin{lemma}[Comparison]
\label{lem:lapcompare}
Let $\Delta(\varphi)=\int\varphi\,d\{(\tilde P-P)-(\tilde Q-Q)\}$, $\widehat\Delta=\Delta(\widehat\varphi)$,
$\Delta^\star=\Delta(\varphi^\star)$, $\widehat T=\mathrm{err}(\widehat\varphi)$, $y=V(\epsilon)+s^\star\Delta^\star$ and
$e=\E_P\widehat\rp-\E_Q\widehat\rp-\alpha=\epsilon-\widehat T$. Then:
\begin{enumerate}[label=(\roman*),leftmargin=*]
\item $e+\widehat\Delta=\bar{\widehat D}\widehat N$, $\bar{\widehat D}\ge y$, and $\widehat Ny+\widehat T\le\epsilon+\widehat\Delta$;
\item if $\kappa_0\epsilon-|\Delta^\star|-\widehat\Delta>0$, then $e\ge\kappa_0\epsilon-|\Delta^\star|-\widehat\Delta$;
\item if $|\Delta^\star|\le\kappa_0\epsilon/4$, then $y\ge V(\epsilon)/2\ge c_V\epsilon^{5/4}/2$.
\end{enumerate}
\end{lemma}

\begin{proof}
Part (i) collects the basic relations. The identity is Lemma~\ref{lem:lapclip}, as
$\tilde L(\widehat\varphi)-L(\widehat\varphi)=\widehat\Delta$. By Lemma~\ref{lem:lapratio} for $\tilde P,\tilde Q$ and
$\varphi^\star$, and \eqref{eq:lappop}, $\bar{\widehat D}\ge s^\star\{L(\varphi^\star)-\alpha+\Delta^\star\}=y$, and
multiplying by $\widehat N>0$ gives the last inequality, since $\bar{\widehat D}\widehat N=e+\widehat\Delta=\epsilon-\widehat
T+\widehat\Delta$.

Parts (ii) and (iii) compare the path at two levels. A fitted ramp with a small excess $e$ has error
$\widehat T=\epsilon-e$ close to $\epsilon$, so at the level $\alpha^\star-2\epsilon$ it still has an excess of about
$\epsilon$. Lemma~\ref{lem:lapratio} at that level then forces its norm $\widehat N$ up, and by (i) a large $\widehat N$
forces $e$ up. The population enters through Proposition~\ref{prop:excess}(i) at $\alpha'=\alpha^\star-2\epsilon$, by
which the population gap is a supporting slope of the path,
\[
s^\star\epsilon\ \le\ V(2\epsilon)-V(\epsilon)\qquad(2\epsilon\le\alpha^\star).
\]
Suppose first $\epsilon<\epsilon_0/2$, so that $V(\epsilon)/V(2\epsilon)\ge\rho_0\ge\kappa_0$.
If $|\Delta^\star|\ge\epsilon V(\epsilon)/V(2\epsilon)$, then $\kappa_0\epsilon-|\Delta^\star|-\widehat\Delta\le-\widehat\Delta<e$
by (i), as $\bar{\widehat D}\widehat N>0$, and (ii) holds; so assume $|\Delta^\star|<\epsilon V(\epsilon)/V(2\epsilon)$. Then
$s^\star|\Delta^\star|<\{V(2\epsilon)-V(\epsilon)\}V(\epsilon)/V(2\epsilon)$, hence $0<y<V(2\epsilon)$ and
$r=y/V(2\epsilon)\in(0,1)$. If $e+\epsilon\le0$, then $\widehat\Delta>\epsilon$ by (i), as $\bar{\widehat D}\widehat N>0$,
contradicting the hypothesis of (ii). Otherwise $2\epsilon-\widehat T=e+\epsilon>0$, Lemma~\ref{lem:lapratio} at the
level $\alpha^\star-2\epsilon$ gives $\widehat N\ge(e+\epsilon)/V(2\epsilon)$, and with (i)
\[
e+\widehat\Delta=\bar{\widehat D}\widehat N\ \ge\ y\widehat N\ \ge\ r(e+\epsilon),
\qquad\text{that is,}\qquad(1-r)e\ \ge\ r\epsilon-\widehat\Delta .
\]
As $s^\star\epsilon\le V(2\epsilon)$, $r\epsilon\ge\epsilon V(\epsilon)/V(2\epsilon)-|\Delta^\star|$, so
$(1-r)e\ge\kappa_0\epsilon-|\Delta^\star|-\widehat\Delta>0$, whence $e>0$ and $e\ge(1-r)e$. For (iii),
$s^\star|\Delta^\star|\le\kappa_0V(2\epsilon)/4\le V(\epsilon)/4$, as $s^\star\le V(2\epsilon)/\epsilon$.

If $\epsilon\ge\epsilon_0/2$, the path is bounded below by a constant: $V(\epsilon)\ge(1-\alpha^\star)c_1(\epsilon_0/2)^{5/4}\ge
2\kappa_0\alpha^\star\ge2\kappa_0\epsilon$, and $s^\star\le2$. For (iii), $s^\star|\Delta^\star|\le\kappa_0\epsilon/2\le
V(\epsilon)/4$. For (ii), if $y\ge0$, (i) and $\widehat N\ge\tfrac12$ give $e+\widehat\Delta\ge y/2\ge
V(\epsilon)/2-|\Delta^\star|$, which is enough; if $y<0$, then $|\Delta^\star|>V(\epsilon)/2\ge\kappa_0\epsilon$, while
$e+\widehat\Delta=\bar{\widehat D}\widehat N>0$, so $e>-\widehat\Delta>\kappa_0\epsilon-|\Delta^\star|-\widehat\Delta$.
\end{proof}

We bound $\Delta^\star$ and $\widehat\Delta$ next, the second for a function fitted on the same data. We do it for
all functions of oscillation at most one at once, in terms of their norm and error, on one event that does not depend
on the level.

\begin{lemma}[Intervals]
\label{lem:lapintervals}
Let $\tilde P$ be the empirical law of $n$ independent draws from $P$, which has a continuous distribution function,
let $\eta'\in(0,1)$, $x=\log\{2(n+1)^2/\eta'\}$ and $c_x=2x+4$. With probability at least $1-\eta'$, every interval
$I\subseteq\R$ satisfies $|\tilde P(I)-P(I)|\le\{2P(I)x/n\}^{1/2}+c_x/n$.
\end{lemma}

\begin{proof}
After the probability integral transform, a grid of at most $(n+1)^2$ intervals brackets every interval to within
$2/n$, so Bernstein's inequality and a union bound over the grid suffice. Almost surely no observation lies in an interval of
constancy of the distribution function $F$ or at one of its endpoints, so $X_i\in I$ if and only if $F(X_i)\in J$ for the interval $J$ with endpoints $F(\inf I),F(\sup I)$ of the
same types, and $P(I)=|J|$; it suffices to treat uniform variables. For each of the at most $(n+1)^2$ intervals $G$ with
endpoints in $\{0,1/n,\dots,1\}$, Bernstein's inequality \citep[Section~2.8]{boucheron2013concentration} for the binomial
count, with deviation $\tau=(2n|G|x)^{1/2}+2x/3$, which satisfies $\tau^2\ge2x(n|G|+\tau/3)$, bounds
$|\tilde P(G)-|G||$ by $(2|G|x/n)^{1/2}+2x/(3n)$ outside an event of probability $2e^{-x}$; a union bound covers all of
them. A general $J$ lies between grid intervals $G^-\subseteq J\subseteq G^+$ with $|G^+|-|J|\le2/n$ and $|J|-|G^-|\le2/n$,
which costs at most $(2\sqrt x+2)/n$; and $2\sqrt x+2x/3+2\le c_x$ as $x>2$.
\end{proof}

\begin{lemma}[Localized deviation]
\label{lem:lapbv}
On the event of Lemma~\ref{lem:lapintervals}, for every bounded interval $I$ and every $h:\R\to[0,1]$ continuous on
the closure of $I$ with total variation $\mathrm V_I(h)$ there,
\[
\Big|\int_Ih\,d(\tilde P-P)\Big|\le\Big[2\{\mathrm V_I(h)+1\}\Big(\int_Ih\,dP\Big)\frac xn\Big]^{1/2}+\{\mathrm V_I(h)+1\}\frac{c_x}n .
\]
\end{lemma}

\begin{proof}
A function of bounded variation is an average of indicators of its upper level sets, each a union of few intervals,
to which Lemma~\ref{lem:lapintervals} applies. By Tonelli's theorem, $\int_Ih\,d(\tilde P-P)=\int_0^1(\tilde P-P)(A_s)\,ds$ and $\int_Ih\,dP=\int_0^1P(A_s)\,ds$, with
$A_s=\{z\in I:h(z)>s\}$. Each component $C$ of the relatively open set $A_s$ other than the rightmost one has
$\sup C\in I\setminus A_s$ and $h(\sup C)=s$, so $A_s$ has at most $N(s)+1$ components, $N(s)$ being the number of
points of $\bar I$ where $h=s$. By Banach's indicatrix theorem \citep{banach1925lignes}, $N$ is measurable and
$\int N(s)\,ds=\mathrm V_I(h)$. Where $N(s)<\infty$, the event and the Cauchy--Schwarz inequality over the components give
$|(\tilde P-P)(A_s)|\le[2\{N(s)+1\}P(A_s)x/n]^{1/2}+\{N(s)+1\}c_x/n$; integrate over $s$ and use the Cauchy--Schwarz
inequality once more.
\end{proof}

\begin{lemma}[Deviation of functions of $\Hh$]
\label{lem:lapdev}
Let $\tilde P,\tilde Q$ be empirical laws of $n$ draws from $P$ and $m$ draws from $Q$ on whose events of
Lemma~\ref{lem:lapintervals} the same $x\ge\log8$ holds, let $K\ge1$, $\pi_K=(P+Q)(\{|z-x_0|>K\})$,
$s_{n,m}=(x/n)^{1/2}+(x/m)^{1/2}$ and $r_{n,m}=1/n+1/m$. Let $\varphi\in\Hh$ with $\operatorname{osc}(\varphi)\le1$,
$N=\max(\|\varphi\|_\Hh,\tfrac12)$ and $T=\mathrm{err}(\varphi)$.
\begin{enumerate}[label=(\alph*),leftmargin=*]
\item $|\Delta(\varphi)|\le a'+\varpi'N^{1/2}+bN$, with $a'=5s_{n,m}+4c_xr_{n,m}+2\pi_K$, $\varpi'=2K^{1/4}s_{n,m}$
and $b=6K^{1/2}c_xr_{n,m}$.
\item Under \eqref{eq:condR}, $|\Delta(\varphi)|\le a+\varpi N^{1/2}T^{3/10}+bN$, with $a=A_0s_{n,m}+7c_xr_{n,m}+2\pi_K$,
$\varpi=A_1K^{1/4}s_{n,m}$, $A_0=4+2(2\bar p)^{1/2}+4(c_\ast r_\ast)^{-1/2}$ and
$A_1=2\{(2\sqrt2\,\bar p)^{1/2}+2^{1/5}(2\sqrt2/(c_\ast r_\ast))^{1/2}\}$.
\end{enumerate}
\end{lemma}

\begin{proof}
We bound $\int\varphi\,d(\tilde P-P)$; the bound for $\tilde Q$ is the same, and the two add. On the window
$W=[x_0-K,x_0+K]$ Lemma~\ref{lem:lapbv} applies, with the variation of $\varphi$ controlled by its norm, and off $W$ only
the tail mass enters. Let $M_\varphi=\sup\varphi$ and $m_\varphi=\inf\varphi$. By \eqref{eq:laphone} the variation of $\varphi$ on
an interval of length $\lambda$ is at most $(2\lambda)^{1/2}N$, and off $W$, as $|\varphi|\le1$,
$|\int_{W^c}\varphi\,d(\tilde P-P)|\le\tilde P(W^c)+P(W^c)\le2\pi^P_K+2(\pi^P_Kx/n)^{1/2}+2c_x/n$, with
$\pi^P_K=P(W^c)$, by the event of Lemma~\ref{lem:lapintervals} for the two half-lines. For (a), write $\varphi=m_\varphi+h$ on $W$ and apply
Lemma~\ref{lem:lapbv} with $\mathrm V_W(h)\le2K^{1/2}N$ and $\int_Wh\,dP\le1$, and the event to $m_\varphi(\tilde P-P)(W)$.

For (b), the error $T$ controls the mass that $\varphi$ carries away from its extreme values where $p\neq q$, and
\eqref{eq:condR} turns this into mass under $P$ away from $x_0$. Note that $T>0$: if
$\operatorname{osc}(\varphi)<1$ by \eqref{eq:laperr}; otherwise $T=0$ would force $\varphi=M_\varphi$ almost everywhere on
$\{p>q\}$ and $\varphi=m_\varphi$ on $\{q>p\}$, which by \eqref{eq:condR} contain punctured neighborhoods of $x_0$ on the
right and on the left, and continuity at $x_0$ would give $M_\varphi=m_\varphi$. Write $\varphi=M_\varphi-h$ on $R_+=(x_0,\infty)$ and
$\varphi=m_\varphi+h$ on $R_-=(-\infty,x_0]$, with $h\in[0,1]$. As $(\tilde P-P)(\R)=0$,
\[
\int\varphi\,d(\tilde P-P)=(M_\varphi-m_\varphi)(\tilde P-P)(R_+)-\int_{R_+}h\,d(\tilde P-P)+\int_{R_-}h\,d(\tilde P-P),
\]
and by \eqref{eq:laperr}, $\int_{R_+}h(p-q)\le T$ and $\int_{R_-}h(q-p)\le T$. Put $w=\min(r_\ast,T^{2/5})$ and split $R_+$
into $I_1=(x_0,x_0+w]$, $I_2=(x_0+w,x_0+K]$ and $(x_0+K,\infty)$. On the short interval $I_1$,
$\mathrm V_{I_1}(h)\le(2w)^{1/2}N$ and $\int_{I_1}h\,dP\le\bar pw$. On $I_2$ the densities differ, \eqref{eq:condR} gives
$p\le(p-q)/(c_\ast w)$, so $\int_{I_2}h\,dP\le T/(c_\ast w)$, while $\mathrm V_{I_2}(h)\le(2K)^{1/2}N$.
Lemma~\ref{lem:lapbv} bounds the two pieces by
\begin{multline*}
(2\sqrt2\,\bar p)^{1/2}w^{3/4}N^{1/2}\Big(\frac xn\Big)^{1/2}+(2\bar p)^{1/2}\Big(\frac xn\Big)^{1/2}
+(\sqrt2N+\sqrt2K^{1/2}N+2)\frac{c_x}n\\
+\Big\{\Big(\frac{2\sqrt2}{c_\ast}\Big)^{1/2}K^{1/4}N^{1/2}+\Big(\frac2{c_\ast}\Big)^{1/2}\Big\}\Big(\frac Tw\Big)^{1/2}\Big(\frac xn\Big)^{1/2} .
\end{multline*}
The width $w$ balances the first and the last term: it gives $w^{3/4}\le T^{3/10}$ and
$(T/w)^{1/2}\le2^{1/5}r_\ast^{-1/2}T^{3/10}$, since if $w=T^{2/5}$ the latter is $T^{3/10}$, and if $w=r_\ast<T^{2/5}$ it is
$T^{3/10}T^{1/5}r_\ast^{-1/2}$ with $T\le2$; also $(T/w)^{1/2}\le(2/r_\ast)^{1/2}$. The half-line $R_-$ is split in the same
way, with $q-p$ for $p-q$. Adding the bounds for the two half-lines, for the term $(M_\varphi-m_\varphi)(\tilde
P-P)(R_+)$, which the event of Lemma~\ref{lem:lapintervals} controls, and off $W$, and using
$2\sqrt2(1+K^{1/2})\le6K^{1/2}$, we get
\[
\Big|\int\varphi\,d(\tilde P-P)\Big|\le A_0\Big(\frac xn\Big)^{1/2}+\frac{7c_x}n+2\pi^P_K
+A_1K^{1/4}N^{1/2}T^{3/10}\Big(\frac xn\Big)^{1/2}+\frac{6K^{1/2}Nc_x}n ,
\]
where $A_0$ collects the constant terms, $2(2\bar p)^{1/2}+4(c_\ast r_\ast)^{-1/2}$ from the two half-lines and at most
$4$ from $(\tilde P-P)(R_+)$ and the tails. For $\tilde Q$, use $q\le(p-q)/(c_\ast w)$ on $I_2$ and $Q(I_1)\le\bar pw$.
\end{proof}

\begin{lemma}[Other cut points]
\label{lem:lapown}
Let $\widehat u\le\tilde u<\tilde t\le\widehat t$, $a_o=(\tilde u-\widehat u)/\widehat s$ and $b_o=(\widehat t-\tilde t)/\widehat
s$, so that $a_o+b_o<1$. The ramp $\tilde\rp$ of $\widehat f$ between $\tilde u$ and $\tilde t$ satisfies
$\E_P\tilde\rp-\E_Q\tilde\rp\ge\E_P\widehat\rp-\E_Q\widehat\rp-\max(a_o,b_o)\,\mathrm{err}(\widehat\rp)/(1-a_o-b_o)$. By
Lemma~\ref{lem:lapclip}, own-term cut points at genuine quantiles have $a_o\le1/(m'\bar{\widehat D}\widehat s)$ and
$b_o\le1/(n'\bar{\widehat D}\widehat s)$.
\end{lemma}

\begin{proof}
Moving the cut points inward rescales the ramp, and the error of the ramp controls what the rescaling costs. As $\widehat f=\widehat u+\widehat s\widehat\rp$, $\tilde\rp=\chi(\widehat\rp)$ with
$\chi(s)=\min[1,\max\{0,(s-a_o)/(1-a_o-b_o)\}]$, and checking the three pieces $[0,a_o]$, $[a_o,1-b_o]$, $[1-b_o,1]$
gives $-a_o(1-s)\le(1-a_o-b_o)\{\chi(s)-s\}\le b_os$ on $[0,1]$. Integrate against $(p-q)_+$ and $(q-p)_+$ and use
$\mathrm{err}(\rp)=\int(1-\rp)(p-q)_++\int\rp(q-p)_+$ for $[0,1]$-valued $\rp$.
\end{proof}

\begin{proposition}[The fitted ramp keeps a linear excess]
\label{prop:lapfitted}
On the events of Lemma~\ref{lem:lapintervals} for $\tilde P$ and $\tilde Q$ with a common $x\ge\log8$, and with the
constants of Lemma~\ref{lem:lapdev} for the sizes $n',m'$, put, in case (i) of Theorem~\ref{thm:laplacerate},
\[
\epsilon_{\mathrm{fit}}=\max\Big\{\frac{24a}{\kappa_0},\frac{108\varpi}{\kappa_0c_V^{1/2}},\Big(\frac{300b}{\kappa_0c_V}\Big)^{4/5},
\Big(\frac2{c_V}\Big)^{20/33}\varpi^{40/33},\Big(\frac{1920}{\kappa_0c_V^2(n'\wedge m')}\Big)^{2/3}\Big\},
\]
and in case (ii)
\[
\epsilon_{\mathrm{fit}}=\max\{24a'/\kappa_0,(60\varpi'\kappa_0^{-1}c_V^{-1/2})^{8/9},(156b\kappa_0^{-1}c_V^{-1})^{4/5},
(576\kappa_0^{-1}c_V^{-2}(n'\wedge m')^{-1})^{2/3}\}.
\]
At every level $\alpha=\alpha^\star-\epsilon\in(0,\alpha^\star)$ with
$\epsilon\ge\epsilon_{\mathrm{fit}}$, the ramp fitted on $\tilde P,\tilde Q$ with genuine-quantile cut points, plain or
own-term, has cut points that do not cross and $\E_P\rp-\E_Q\rp-\alpha\ge\kappa_0\epsilon/2$.
\end{proposition}

\begin{proof}
In case (i) we bound the fluctuation of the population ramp, then the norm and the error of the fitted ramp by a
fixed-point argument, and then its fluctuation; the constants $24$, $108$ and $300$ in $\epsilon_{\mathrm{fit}}$ are
chosen so that $a$, $\varpi$ and $b$ each contribute at most $\kappa_0\epsilon/12$ at the end; the fourth term gives $\varpi^5y^{-5/2}\le\epsilon$, and the fifth, with $1920=48\cdot40$, bounds the shifts $a_o,b_o$ of the own-term cut points by $\kappa_0/40$. As
$\|\varphi^\star\|_\Hh\le c_V^{-1}\epsilon^{-1/4}$ and $\mathrm{err}(\varphi^\star)\le\epsilon$, Lemma~\ref{lem:lapdev}(b) and
the first three terms of $\epsilon_{\mathrm{fit}}$ give
\[
|\Delta^\star|\le a+c_V^{-1/2}\varpi\epsilon^{7/40}+c_V^{-1}b\epsilon^{-1/4}\le\kappa_0\epsilon\Big(\frac1{24}+\frac1{108}+\frac1{300}\Big)<\frac{\kappa_0\epsilon}4 .
\]
By Lemma~\ref{lem:lapcompare}(iii), $y\ge c_V\epsilon^{5/4}/2$, so $b\le y/4$. The fitted ramp now appears on both sides:
Lemma~\ref{lem:lapcompare}(i) and Lemma~\ref{lem:lapdev}(b) for $\widehat\varphi$ give
$\widehat Ny+\widehat T\le\epsilon+a+\varpi\widehat N^{1/2}\widehat T^{3/10}+b\widehat N$. The weighted arithmetic--geometric mean inequality $u^{1/2}v^{3/10}w^{1/5}\le u/2+3v/10+w/5$ with $u=\widehat Ny$,
$v=\widehat T$ and $w=\varpi^5y^{-5/2}$ separates the unknowns on the right, and with $b\widehat N\le\widehat Ny/4$ gives
\[
\frac{\widehat Ny}4+\frac{7\widehat T}{10}\le\epsilon+a+\frac{\varpi^5y^{-5/2}}5\le3\epsilon ,
\]
the fourth term of $\epsilon_{\mathrm{fit}}$ giving $\varpi^5y^{-5/2}\le\epsilon$. Hence $\widehat N\le24c_V^{-1}\epsilon^{-1/4}$
and $\widehat T\le5\epsilon$: up to constants, the fitted ramp has the norm and the error of the population one.
Lemma~\ref{lem:lapdev}(b) once more gives $\widehat\Delta\le a+8c_V^{-1/2}\varpi+24c_V^{-1}b\epsilon^{-1/4}$, so that
$|\Delta^\star|+\widehat\Delta\le2a+9c_V^{-1/2}\varpi+25c_V^{-1}b\epsilon^{-1/4}\le\kappa_0\epsilon/4$, and by Lemma~\ref{lem:lapcompare}(ii) the plain ramp has excess at least $3\kappa_0\epsilon/4$. For own-term cut points,
$\bar{\widehat D}\widehat s\ge y/\widehat N\ge c_V^2\epsilon^{3/2}/48$, so by Lemma~\ref{lem:lapown} and the fifth term
$\max(a_o,b_o)\le\kappa_0/40$, and the loss is at most $(\kappa_0/40)(20/19)\,5\epsilon<\kappa_0\epsilon/4$.

Case (ii) is the same argument with part (a) of Lemma~\ref{lem:lapdev}, which has no factor $T^{3/10}$. It gives
$|\Delta^\star|\le a'+c_V^{-1/2}\varpi'\epsilon^{-1/8}+c_V^{-1}b\epsilon^{-1/4}\le\kappa_0\epsilon/4$, so
$y\ge c_V\epsilon^{5/4}/2\ge4b$. With $\varpi'\widehat N^{1/2}\le\widehat Ny/4+\varpi'^2/y$ and $\varpi'^2/y\le\epsilon$
(second term), $\widehat Ny/2+\widehat T\le3\epsilon$, so $\widehat N\le12c_V^{-1}\epsilon^{-1/4}$, $\widehat T\le3\epsilon$ and
$|\Delta^\star|+\widehat\Delta\le2a'+5c_V^{-1/2}\varpi'\epsilon^{-1/8}+13c_V^{-1}b\epsilon^{-1/4}\le\kappa_0\epsilon/4$; the own
terms have $\max(a_o,b_o)\le24/\{c_V^2(n'\wedge m')\epsilon^{3/2}\}\le\kappa_0/24$ and cost less than $\kappa_0\epsilon/4$.
\end{proof}


\begin{proof}[Proof of Theorem~\ref{thm:laplaceconst}]
The structure is that of the proof of Theorem~\ref{thm:lorate}, with Proposition~\ref{prop:lapfitted} in place of the
margin: one event makes every fitted ramp good, and a union over the levels handles the bets.

We first fix the fitting event. Take
\[
x=\log\{16(n\vee m+1)^2/\eta\},\qquad K=\max\{1,C'\log(n\vee m)\},
\]
so that the window of Lemma~\ref{lem:lapdev} leaves a tail mass of order $1/(n\vee m)$. Let $\mathcal E_{\mathrm{fit}}$
be the intersection of the events of Lemma~\ref{lem:lapintervals}, each with $\eta'=\eta/8$, for the $X$'s and the
$Y$'s of batch $0$ and of batches $0$ and $1$. These events weaken as $x$ grows, so all four hold with this $x$, and
$\Pp(\mathcal E_{\mathrm{fit}}^c)\le\eta/2$. The constants of Lemma~\ref{lem:lapdev} and the last term of
$\epsilon_{\mathrm{fit}}$ decrease in the sample sizes, so the $\epsilon_{\mathrm{fit}}$ of batch $0$ serves for both
fits.

Next we choose the distance $\epsilon$ from $\alpha^\star$. Let
\[
\widetilde T_{\mathrm{bet}}(J')=\Big(\frac2N\Big)^{1/2}\Big\{\Big(1+\frac{16}{v^2_{\min}}\Big)L_\delta^{1/2}
+\frac2{v_{\min}}\big(\log(2J'/\eta)\big)^{1/2}\Big\},
\]
the $T_{\mathrm{bet}}(\eta,J')$ of \eqref{eq:loconst} divided by $2\Delta_\phi$. With $J_\epsilon$ and $a_\epsilon$ as in
Theorem~\ref{thm:lorate}, fix $\epsilon\in[\epsilon_{\mathrm{fit}},\alpha^\star)$ with
$\kappa_0\epsilon\ge2\widetilde T_{\mathrm{bet}}(J_\epsilon)$, and put $\varrho_\ast=\kappa_0\epsilon/2$. Write each grid
level $\alpha_j\le\alpha^\star-\epsilon$ as $\alpha_j=\alpha^\star-\epsilon\xi_j$, so that $\xi_j\ge1$ and
$\epsilon\xi_j\ge\epsilon_{\mathrm{fit}}$. On $\mathcal E_{\mathrm{fit}}$, Proposition~\ref{prop:lapfitted} at the level
$\alpha_j$ shows that both batches have a ramp with
\[
\E_P\rp_b-\E_Q\rp_b-\alpha_j\ \ge\ \varrho_\ast\xi_j .
\]

The bets then win at every such level. As $1>\kappa_0\epsilon\ge10(2L_\delta/N)^{1/2}$, we have $N\ge200L_\delta$, and
the first bound of \eqref{eq:betbounds}, with $\mathcal E=\mathcal E_{\mathrm{fit}}$ and $\varrho=\varrho_\ast\xi_j$,
gives
\[
\Pp\big(\mathcal E_{\mathrm{fit}}\cap\{H_{0,\alpha_j}\text{ is not rejected}\}\big)
\le\exp\Big\{-\frac{v^2_{\min}}{4L_\delta}\big(\varrho_\ast\xi_jS-C_0\big)_+^2\Big\},
\]
with $S=(NL_\delta/2)^{1/2}$ and $C_0=L_\delta(1+16/v^2_{\min})$. Multiplied by $S$, the condition
$\varrho_\ast\ge\widetilde T_{\mathrm{bet}}(J_\epsilon)$ reads
$a_\ast=\varrho_\ast S-C_0\ge2v_{\min}^{-1}\{L_\delta\log(2J_\epsilon/\eta)\}^{1/2}$. Since
$\varrho_\ast\xi_jS-C_0\ge a_\ast\xi_j$, each term is at most $e^{-A\xi_j^2}$ with
$A=v^2_{\min}a_\ast^2/(4L_\delta)\ge\log(2J_\epsilon/\eta)>\log2$. The summation over the levels in the proof of
Theorem~\ref{thm:lorate}, with $\gamma=1$, bounds the sum by $\eta/2$, and \eqref{eq:downclosed} gives
$\Pp(\widehat\alpha<a_\epsilon)\le\eta$. Put $\epsilon_{n,m}=\max\{\epsilon_{\mathrm{fit}},2\widetilde
T_{\mathrm{bet}}(J)/\kappa_0\}$, which is either at least $\alpha^\star$ or satisfies the requirements on $\epsilon$.
With Corollary~\ref{cor:lower}, with probability at least $1-\delta-\eta$,
\[
\alpha^\star-\Delta_{\mathcal G}-\min\{\epsilon_{n,m},\alpha^\star\}\ \le\ \widehat\alpha\ \le\ \alpha^\star .
\]

It remains to read the orders of the terms of $\epsilon_{\mathrm{fit}}$ and of $\widetilde T_{\mathrm{bet}}$. As in the
proof of Theorem~\ref{thm:lorate}(iii), $n'\wedge m'\ge(n\wedge m)/5$ and $N\ge(n\wedge m)/2$ for $n\wedge m\ge10$, so
that
\[
s_{n',m'}\le2\Big(\frac{5x}{n\wedge m}\Big)^{1/2},\qquad r_{n',m'}\le\frac{10}{n\wedge m},\qquad
\pi_K\le\frac{C''}{n\vee m},
\]
with $C''$ depending on $C'$ and, in case (i), on $c_\ast r_\ast$, and
$K^{1/4}x^{1/2}\le C'''(\log(n\vee m))^{1/4}L_{n,m}^{1/2}$ with $C'''\ge1$ depending only on $C'$. If
$K^{1/2}x>n\wedge m$, the bounds of the theorem exceed $1>\alpha^\star$ once $C\ge C'''$, and there is nothing to prove.
Otherwise $u=K^{1/4}\{x/(n\wedge m)\}^{1/2}\le1$, and $u^{8/5}=\{K^{1/2}x/(n\wedge m)\}^{4/5}$, $u^{40/33}$,
$(n\wedge m)^{-2/3}$ and $x/(n\wedge m)$, the other quantities through which $a$, $\varpi$, $b$ and the last term enter
$\epsilon_{\mathrm{fit}}$ in case (i), are all at most $u$. Hence
\[
\epsilon_{\mathrm{fit}}\le Cu\ \text{ in case (i)},\qquad
\epsilon_{\mathrm{fit}}\le C\Big\{\frac{K^{1/2}x}{n\wedge m}\Big\}^{4/9}\ \text{ in case (ii)},
\]
and, as $v^2_{\min}\le4$,
\[
\widetilde T_{\mathrm{bet}}(J)\le\frac{2\{20L_\delta^{1/2}+4(\log(2J/\eta))^{1/2}\}}{v^2_{\min}(n\wedge m)^{1/2}}.
\]
\end{proof}

Near every pair covered by part (i) no valid lower bound does better than $(n\wedge m)^{-1/2}$. The pairs of Proposition~\ref{prop:lower2pt} are separated
and do not satisfy these hypotheses, so we perturb $P$ away from $x_0$ instead.

\begin{lemma}[A local two-point bound]
\label{lem:laplocal}
Let $(P,Q)$ satisfy the hypotheses of Theorem~\ref{thm:laplacerate}(i), let $r\in(0,r_\ast]$ be such that $p,q$ are continuous
and positive on $[x_0-2r,x_0+2r]$, put $A_1=(x_0+r,x_0+2r)$, $B_1=(x_0-2r,x_0-r)$, $\pi=\min\{P(A_1),P(B_1)\}>0$ and, for
$h\in[0,1]$,
\[
P_h=P+h\pi\Big\{\frac{P(\cdot\cap A_1)}{P(A_1)}-\frac{P(\cdot\cap B_1)}{P(B_1)}\Big\}.
\]
Then $P_h$ is a probability measure, $\TV(P_h,Q)=\alpha^\star+h\pi$, and $(P_h,Q)$ satisfies the hypotheses of
Theorem~\ref{thm:laplacerate}(i), with $c_\ast/2$, $2\bar p$ and a tail constant depending only on $C'$ and $r$. Every
$\tilde\alpha$, a measurable function of the samples and of an independent randomization, with
$\Pp_{P,Q}(\tilde\alpha>\alpha^\star)\le\delta$ satisfies
$\Pp_{P_h,Q}\{\TV(P_h,Q)-\tilde\alpha\ge h\pi\}\ge1-\delta-h(n\pi)^{1/2}$. The same holds for the second sample, with $Q$
moved from $A_1$ to $B_1$, $m$ in place of $n$ and $\pi'=\min\{Q(A_1),Q(B_1)\}$ in place of $\pi$.
\end{lemma}

\begin{proof}
Moving mass $h\pi$ from $B_1$ to $A_1$, away from $x_0$, raises the total variation by $h\pi$ without touching the
crossing, and costs little in $\chi^2$-divergence. The density of $P_h$ is $p\{1+h\pi/P(A_1)\}$ on $A_1$, $p\{1-h\pi/P(B_1)\}\ge0$ on $B_1$ and $p$ elsewhere; it jumps
at $x_0\pm r$ and $x_0\pm2r$, which is harmless, as the hypotheses ask continuity near $x_0$ only. By
\eqref{eq:condR}, $p>q$ almost everywhere on $A_1$ and $q>p$ on $B_1$, so $p_h-q$ has the sign of $p-q$ everywhere and
$\TV(P_h,Q)=\int_{\{p>q\}}(p_h-q)=\alpha^\star+h\pi$. Near $x_0$ nothing changes, so the hypotheses of
Proposition~\ref{prop:newsobolev} hold; on $A_1$, $|p_h-q|\ge|p-q|$ and $p_h+q\le2(p+q)$, and on $B_1$, $|p_h-q|\ge|p-q|$ and
$p_h\le p$, which gives \eqref{eq:condR} with $c_\ast/2$. Further, $\chi^2(P_h\|P)=(h\pi)^2\{1/P(A_1)+1/P(B_1)\}\le2h^2\pi$, so,
by $\mathrm{KL}\le\chi^2$ for each of the $n$ observations and Pinsker's inequality, the laws of the data and the randomization under $(P,Q)$ and
$(P_h,Q)$ are at total variation at most $\{n\chi^2(P_h\|P)/2\}^{1/2}\le h(n\pi)^{1/2}$, and on $\{\tilde\alpha\le\alpha^\star\}$ the
error at $(P_h,Q)$ is at least $h\pi$. For the second sample, $Q_h=Q+h\pi'\{Q(\cdot\cap B_1)/Q(B_1)-Q(\cdot\cap A_1)/Q(A_1)\}$
with $\pi'=\min\{Q(A_1),Q(B_1)\}$ gives $\TV(P,Q_h)=\alpha^\star+h\pi'$ in the same way.
\end{proof}

With $h=(1-\delta)/\{2(n\pi)^{1/2}\}$, which is at most $1$ once $n\pi\ge1/4$, a valid lower bound falls short at $(P_h,Q)$
by at least $\tfrac12(1-\delta)(\pi/n)^{1/2}$ with probability at least $(1-\delta)/2$. The constants of
Proposition~\ref{prop:newsobolev} for $(P_h,Q)$ exist; whether they stay bounded uniformly in $h$, as a minimax
statement over a class containing these pairs would need, is open.

\subsection{The two-point lower bound and the minimax rate}
\label{sec:optproofs}

This subsection proves Proposition~\ref{prop:lower2pt} and Corollary~\ref{cor:minimax} of the main text. The
perturbed pairs of Proposition~\ref{prop:lower2pt} satisfy the margin condition with a constant that does not depend on
$a$ or $h$.

\begin{lemma}[The two-point pairs satisfy the margin]
\label{lem:twopointmargin}
Let (A1) and (A2) of Assumption~\ref{ass:regularity} hold, let $R$, $P_0$ and $Q_0$ be probability measures
concentrated on pairwise disjoint measurable sets $S_R$, $S_P$ and $S_Q$, define $P_a$, $Q_a$, $P_{a,h}$ and $Q_{a,h}$
as in Proposition~\ref{prop:lower2pt}, and suppose some $f\in\Hh$ takes the value $0$ on $S_R$, $1$ on $S_P$ and $-1$
on $S_Q$. Then for every $a\in(0,1)$ and $h\in[0,1-a)$, each of the pairs $(P_{a,h},Q_a)$ and $(P_a,Q_{a,h})$, with
$\alpha^\star=a+h$, satisfies $D_{\alpha^\star-\epsilon}\ge\epsilon/\|f\|_\Hh$ for all $\epsilon\in[0,\alpha^\star]$,
that is, Assumption~\ref{ass:margin} with $\gamma=1$ and $c_0=1/\|f\|_\Hh$.
\end{lemma}

\begin{proof}
Consider $(P_{a,h},Q_a)$ and fix $\alpha\in[0,\alpha^\star]$. Under $P_{a,h}$, $f(X)\in\{0,1\}$ almost surely with
$P_{a,h}(f(X)=1)=a+h$. Since $1-\alpha\ge1-a-h$, an $\alpha$-trimming $P'=gP_{a,h}$ puts mass at most
$\beta(1-a-h)\le1$ on $\{f=0\}$, so $\E_{P'}f\ge1-\beta(1-a-h)=(a+h-\alpha)/(1-\alpha)$. Under $Q_a$,
$f(Y)\in\{0,-1\}$ almost surely with $Q_a(f(Y)=-1)=a$, and an $\alpha$-trimming $Q'$ puts mass at most
$\min\{1,\beta(1-a)\}$ on $\{f=0\}$, so $\E_{Q'}f\le-(1-\beta(1-a))_+=-(a-\alpha)_+/(1-\alpha)\le0$. By the reproducing
property and the Cauchy--Schwarz inequality, for all such $P'$ and $Q'$,
\[
\MMD(P',Q')\ge\frac{\E_{P'}f-\E_{Q'}f}{\|f\|_\Hh}\ge\frac{a+h-\alpha}{(1-\alpha)\|f\|_\Hh}\ge\frac{\alpha^\star-\alpha}{\|f\|_\Hh} .
\]
Taking the infimum over $P'$ and $Q'$ and putting $\alpha=\alpha^\star-\epsilon$ gives the claim. The pair
$(P_a,Q_{a,h})$ is the same computation with $-f$ in place of $f$ and the roles of the two laws exchanged.
\end{proof}

\begin{proof}[Proof of Proposition~\ref{prop:lower2pt}]
The two pairs are close in $\chi^2$, so a bound valid at the first must often fall short at the second, by Le Cam's
two-point method. The three measures are mutually singular. Hence $P_{a,h}\wedge Q_a$ has mass $\min(1-a-h,1-a)=1-a-h$ on $S_R$ and
none elsewhere, so $\TV(P_{a,h},Q_a)=a+h$, and likewise $\TV(P_a,Q_{a,h})=a+h$ and $\TV(P_a,Q_a)=a$. The density of
$P_{a,h}$ with respect to $P_a$ is $(1-a-h)/(1-a)$ on $S_R$ and $(a+h)/a$ on $S_P$, so
\[
\chi^2(P_{a,h}\,\|\,P_a)=(1-a)\Big(\frac h{1-a}\Big)^2+a\Big(\frac ha\Big)^2=\frac{h^2}{a(1-a)} .
\]
Let $\mathbb Q_0$ and $\mathbb Q_1$ be the laws of the data and the randomization under $(P_a,Q_a)$ and under
$(P_{a,h},Q_a)$. They differ only in the law of the $n$ observations of the first sample, and the Kullback--Leibler
divergence adds over independent coordinates, so with $\mathrm{KL}\le\log(1+\chi^2)\le\chi^2$,
$\mathrm{KL}(\mathbb Q_1\|\mathbb Q_0)\le nh^2/\{a(1-a)\}$. By Pinsker's inequality
$\TV(\mathbb Q_0,\mathbb Q_1)\le\{\mathrm{KL}(\mathbb Q_1\|\mathbb Q_0)/2\}^{1/2}\le h\{n/(2a(1-a))\}^{1/2}$. Therefore
\[
\Pp_{P_{a,h},Q_a}(\tilde\alpha\le a)\ge\Pp_{P_a,Q_a}(\tilde\alpha\le a)-\TV(\mathbb Q_0,\mathbb Q_1)
\ge1-\delta-h\sqrt{n/\{2a(1-a)\}},
\]
and on $\{\tilde\alpha\le a\}$ the error at $(P_{a,h},Q_a)$, whose total variation is $a+h$, is at least $h$. The same
argument with the second sample perturbed gives the bound with $m$.
\end{proof}

\begin{proof}[Proof of Corollary~\ref{cor:minimax}]
Write $\bar n=n\wedge m$. For (i), let $R$, $P_0$ and $Q_0$ be the point masses at $x_2$, $x_1$ and $x_3$, put $a=a_0$ and
$h=c_{\mathrm{lo}}\bar n^{-1/2}$, so that $a_0+h\le a_1<1$ by the condition on $\bar n$. The pair $(P_{a_0},Q_{a_0})$ belongs to
$\mathcal M$ by Example~\ref{ex:separated}, since $c_0\le1/\|f\|_\Hh\le2/\|f\|_\Hh$, and the two perturbed pairs belong
to $\mathcal M$ by Lemma~\ref{lem:twopointmargin}. If $\bar n=n$, Proposition~\ref{prop:lower2pt} at the pair
$(P_{a_0,h},Q_{a_0})$ gives an error of at least $h$ with probability at least
$1-\delta-h\{n/(2a_0(1-a_0))\}^{1/2}=(1-\delta)/2$; if $\bar n=m$, the pair $(P_{a_0},Q_{a_0,h})$ gives the same with $m$.

For (ii), validity is Corollary~\ref{cor:lower}. On the grid $\{j\bar n^{-1/2}\}$ the smallest gap is $\Delta_{\min}=\bar n^{-1/2}$ and
the mesh below $\alpha^\star$ is at most $\bar n^{-1/2}$. Apply Theorem~\ref{thm:lorate}(i) with $\epsilon=C_\epsilon
\bar n^{-1/2}$: then $J_\epsilon\le2(1+C_\epsilon)$ does not grow with $\bar n$, and by Theorem~\ref{thm:lorate}(iii),
$r_1\le12/\bar n$ and the number of pairs is at least $\bar n/2$, so $\beta_\star\max\{T_{\mathrm{fit}}(\eta),
T_{\mathrm{bet}}(\eta,J_\epsilon)\}\le C''(C_\epsilon)\bar n^{-1/2}$ with $\beta_\star\le(1-a_1)^{-1}$ and $C''(C_\epsilon)$
growing like $\{\log(1+C_\epsilon)\}^{1/2}$. Choose $C_\epsilon$ with $c_0C_\epsilon\ge C''(C_\epsilon)$ and let
$\bar n\ge10$, so that Theorem~\ref{thm:lorate}(iii) applies; every step of the proof holds for these sizes, so $n_1=10$
will do. If $\epsilon>\alpha^\star$, then $\widehat\alpha\ge0>\alpha^\star
-\epsilon$ and there is nothing to prove. Otherwise $\epsilon\in(0,\alpha^\star]$, Theorem~\ref{thm:lorate}(i) gives
$\Pp(\widehat\alpha<a_\epsilon)\le\eta$, and $a_\epsilon\ge\alpha^\star-\epsilon-\bar n^{-1/2}$. In both
cases $\Pp(\alpha^\star-\widehat\alpha>(C_\epsilon+1)\bar n^{-1/2})\le\eta$. Take $C=C_\epsilon+1$.
\end{proof}
